% Canonical static Paper-III manuscript. Numerical authority:
% output/pdf/paper_iii_static_qse_source_anchor_20260829_v2/
% This file is intentionally expansive. Pruning is an authorial pass, not a
% condition of the evidence build.
\documentclass[
  reprint,
  amsmath,
  amssymb,
  aps,
  pra,
  nofootinbib,
  superscriptaddress,
  floatfix
]{revtex4-2}

\usepackage[T1]{fontenc}
\usepackage{lmodern}
\usepackage{graphicx}
\usepackage{booktabs}
\usepackage{bm}
\usepackage{mathtools}
\usepackage{xcolor}
\usepackage{hyperref}
\hypersetup{hidelinks}
\raggedbottom

\newcommand{\ii}{\mathrm{i}}
\newcommand{\ket}[1]{|#1\rangle}
\newcommand{\bra}[1]{\langle #1|}
\newcommand{\braket}[2]{\langle #1|#2\rangle}
\newcommand{\Tr}{\operatorname{Tr}}
\newcommand{\rank}{\operatorname{rank}}
\newcommand{\spanop}{\operatorname{span}}
\newcommand{\cond}{\operatorname{cond}}
\newcommand{\argmin}{\operatorname*{arg\,min}}
\newcommand{\OSQSE}{\mathrm{OS\mbox{-}QSE}}
\newcommand{\DQ}{\mathrm{DQ}}
\newcommand{\RTK}{\mathrm{RTK}}
\newcommand{\mathcalR}{\mathcal{R}}
\newcommand{\score}{\mathcal{S}}
\newcommand{\phmax}{n_{\mathrm{ph,max}}}
\newcommand{\nph}{n_{\mathrm{ph}}}
\newcommand{\sourceanchor}{\textit{static-QSE numerical source anchor}}
\newcommand{\inlinetablecaption}[2]{%
  \refstepcounter{table}\label{#1}%
  \noindent\textsc{Table~\thetable.} #2\par\vspace{0.6ex}%
}
\newcommand{\inlinefigurecaption}[2]{%
  \refstepcounter{figure}\label{#1}%
  \noindent\textsc{Figure~\thefigure.} #2\par\vspace{0.6ex}%
}
\pdfinfoomitdate=1
\pdfsuppressptexinfo=15

\begin{document}

\title{Cost-Weighted Operator-Selected Quantum Subspace Expansion for Static Spectra of Correlated Fermion--Boson Systems}

\author{Jake Strobel}
\author{Jason Simoni}
\affiliation{University of Wisconsin--Madison, Madison, Wisconsin 53706, USA}

\date{August 29, 2026}

\begin{abstract}
We introduce cost-weighted operator-selected quantum subspace expansion
(OS-QSE) for constructing low-energy spectral windows of correlated
fermion--boson Hamiltonians.  Candidate records are evaluated by the decrease
of the full proposed supported-pencil Ky--Fan sum, divided by a clamped,
fixed-pool-normalized graph-span cost, while overlap novelty acts only as an
admissibility floor.
Atomic singleton exchanges maintain the admitted support under recomputed
spectral, cost, retained-rank, and conditioning guards.  Exact-seed noiseless
benchmarks comprise twelve Hubbard--Holstein dimers at two phonon cutoffs,
three Peierls--Hubbard regimes, and a ten-qubit finite-box single-impurity
Anderson model.  OS-QSE reaches a six-root maximum energy error of $10^{-6}$
in all sixteen cases without residual-synthesis continuation.  At each
route's first threshold crossing, its summed two-qubit circuit-structure
annotation is smaller than that of independently initialized direct quantum
Davidson.  Fixed-policy single-seed and six-probe block real-time-Krylov
comparators reach ten and eleven of the sixteen target windows,
respectively; their complementary failures expose initialization-dependent
spectral support and are never combined into a post-hoc oracle route.  A
controlled Hubbard--Holstein pool ablation repairs two six-root coverage
failures at the $10^{-4}$ target when encoded-boson completion records are
added.  Static density,
displacement, double-occupancy, and magnetization spectra independently test
transition structure.  We distinguish online graph-span cost, post-run
circuit-structure annotations, and route-specific logical estimator queries;
finite-shot, total-work, approximate-seed, and time-dynamics claims remain
outside the benchmark.
\end{abstract}

\maketitle

\section{Introduction}
\label{sec:introduction}

Low-energy spectra organize the response of interacting quantum systems.  A
ground-state approximation alone does not identify which excited directions
carry a density peak, a phonon sideband, a bond response, or an impurity
fluctuation.  In mixed fermion--boson problems this difficulty is amplified:
electronic rearrangements, bosonic quanta, and correlation-dressed products
can all be required within the same narrow energy window.  Hubbard and
Holstein models provide canonical electronic and electron--phonon limits
\cite{Hubbard1963,Holstein1959}; Peierls or Su--Schrieffer--Heeger coupling
moves the boson from a site-energy channel to a bond channel
\cite{SuSchriefferHeeger1979}; and finite-bath impurity models supply a
distinct test in which local interaction competes with hybridization and bath
kinetics.  These examples are small enough for exact retrospective audits yet
rich enough to make a fixed response alphabet nontrivial.

Quantum subspace expansion (QSE) replaces preparation of each excited state
by measurement of a projected Hamiltonian and overlap pencil in a family of
operator-generated states \cite{McClean2017QSE,Colless2018QSE}.  The same
matrix viewpoint underlies generalized QSE, virtual subspaces, quantum
Krylov diagonalization, Green-function constructions, and several
equation-of-motion approaches
\cite{Urbanek2020VQSE,Yoshioka2022GQSE,CortesGray2022QKSD,
Motta2024SubspaceReview,Jamet2022GreenQSE,Ollitrault2020qEOM}.  It is most
attractive when a physically structured manifold is much smaller than the
encoded Hilbert space.  Its central design question is consequently not how
to solve a generalized eigenproblem---that part is standard---but which
operator images should be paid for and retained.

The obvious complete-class strategy measures every record in a chosen
response family.  This can miss the target window if the family is not
expressive, and it can be unnecessarily expensive if the family is broad.
Conversely, residual-synthesizing eigensolvers such as quantum Davidson can
generate directions outside a fixed operator pool, but their synthesized
actions may be costly and their discrete growth can overshoot a requested
accuracy \cite{Tkachenko2024QDavidson}.  Quantum Krylov constructions move
the design choice into a source state or source block and a sequence of
Hamiltonian evolutions.  Their formal span can be large while remaining
blind to an eigenvector having zero source overlap.  These routes are not
ordered universally; each exposes a different balance between physical
structure, spectral coverage, and construction cost.

OS-QSE addresses this balance as a cost-weighted spectral-manifold
construction problem.  The ideal objective is the least-cost support that
resolves a specified set of roots.  The online implementation cannot use the
exact roots, so it greedily evaluates the actual supported pencil proposed by
each admissible record and computes its multi-root Ky--Fan gain.  Novelty is
kept as a floor, not a competing utility.  The score is therefore aligned
with the six-root target rather than with a one-root angular residual proxy.
The selected support may subsequently undergo an atomic delete--insert
transaction, but only after the complete patched pencil passes the declared
spectral, cost, rank, and conditioning tests.

This paper makes four bounded contributions.  First, it states the static
spectral-window task and a production selector that uses full proposed-pencil
Ky--Fan gain per cost.  Second, it defines an explicit encoded-boson
completion for the Hubbard--Holstein pool and demonstrates where the
additional directions repair coverage.  Third, it compares OS-QSE with
independently initialized Davidson, complete fixed-class QSE, and two fixed
real-time-Krylov policies under separate structural resource coordinates.
Fourth, it evaluates transition spectra as an independent test of the
selected states.  The numerical evidence supports the proposed method across
the declared sixteen-case benchmark; it does not establish superiority over
all Krylov or variational excited-state methods, lower shot complexity, or
lower end-to-end runtime.

Notation is kept continuous with the companion construction papers
\cite{PaperI,PaperII}: $r\in\mathcal R$ labels requested roots,
$\mathcal B$ is the admitted support, $\Delta$ is the full-pencil proposal gain,
$K$ is the declared construction-cost coordinate, and $\score$ is the
gain-per-cost acquisition score (with the implementation's explicit
normalizations and cost floor).  The delete--insert pair
$(\mathcal B^-,\mathcal B^+)$ likewise retains the Paper-II atomic-patch
meaning.  Roman $S$ is reserved below for the QSE overlap matrix, preventing
the score and metric from being conflated.

The present article is static and frequency-domain.  Real-time propagation,
frozen-QSE dynamics, and McLachlan refinement are intentionally deferred to a
separate study.  Time evolution appears here only as the state-construction
primitive of the two real-time-Krylov comparators.  This separation keeps the
method claim tied to the evidence actually acquired.

\section{Spectral task and benchmark families}
\label{sec:benchmark_task}

\subsection{Target sector and seed convention}

Let $H_0$ be a finite encoded Hamiltonian and let $\mathcal H_{\mathfrak s}$
be a conserved particle-number sector.  We target
\begin{equation}
  \mathcalR=\{1,\ldots,R\},\qquad R=6,
  \label{eq:target_set_intro}
\end{equation}
where $r$ labels the six lowest excited roots in that sector.  Every
production route accepts a prepared-state interface.  For the benchmark,
however, the supplied seed $\ket{\psi_0}$ is the exact ground state of the
same finite Hamiltonian, sector, and bosonic cutoff.  This choice isolates
excited-manifold construction from ground-state preparation.  It is not
evidence for approximate-seed robustness, and its preparation cost is
excluded consistently from the construction comparisons.

Exact excited energies and eigenvectors are separated from the online route.
They are loaded only after a trace is complete, for one-to-one root
assignment where needed, threshold crossings, source-coverage audits, and
plots.  Thus exact ground-state information is an explicit benchmark input,
whereas exact excited-state information is retrospective.

\subsection{Hubbard--Holstein dimer}
\label{sec:hh_model}

The main matrix uses a two-site, spinful Hubbard--Holstein Hamiltonian with
open boundary conditions,
\begin{align}
 H_{\mathrm{HH}}
 =&-t\sum_{\sigma}
 \left(c_{1\sigma}^{\dagger}c_{2\sigma}
 +c_{2\sigma}^{\dagger}c_{1\sigma}\right)
 +U\sum_{i=1}^{2}n_{i\uparrow}n_{i\downarrow}
 \nonumber\\
 &+\omega_0\sum_{i=1}^{2}\left(b_i^{\dagger}b_i+\frac12\right)
 +g\sum_{i=1}^{2}(b_i+b_i^{\dagger})(n_i-1) .
 \label{eq:hh_hamiltonian}
\end{align}
We set $t=\omega_0=1$, use the $(N_\uparrow,N_\downarrow)=(1,1)$
sector, and encode each local phonon register in binary.  The electronic
interaction is $U\in\{0.25,1.25,8\}$ and the electron--phonon coupling is
$g\in\{0.353553390593,0.790569415042\}$.  Their Cartesian product defines
six regimes.  The cutoffs $\phmax=1$ and $3$ require six and eight qubits,
respectively; their fixed-particle sector dimensions are 16 and 64.  The
Cartesian matrix gives twelve benchmark cells.  These are same-cutoff
algorithmic tests; two cutoffs do not constitute a physical convergence
study toward an untruncated oscillator.

\subsection{Peierls--Hubbard dimer}
\label{sec:peierls_model}

The second family places a single encoded bond phonon on the hopping channel,
\begin{align}
 H_{\mathrm P}
 =&-t\,T+gX_bT
 +U\sum_{i=1}^{2}n_{i\uparrow}n_{i\downarrow}
 +\omega_0\left(b^{\dagger}b+\frac12\right),
 \label{eq:peierls_hamiltonian}\\
 T=&\sum_\sigma
 \left(c_{1\sigma}^{\dagger}c_{2\sigma}
 +c_{2\sigma}^{\dagger}c_{1\sigma}\right),
 \qquad X_b=b+b^{\dagger}.
\end{align}
Here $t=\omega_0=1$, $\phmax=3$, and the sector is again $(1,1)$.
The three cases are $(U,g)=(0.25,0.353553390593)$,
$(0.25,0.790569415042)$, and $(8,0.353553390593)$.  The blocked fermion
register plus the two-qubit bond-phonon register uses six qubits; the target
sector has dimension $16$.

\subsection{Finite-box single-impurity Anderson model}
\label{sec:siam_model}

The transfer matrix is completed by the particle--hole-symmetric finite-box
single-impurity Anderson model (SIAM),
\begin{align}
 H_{\mathrm{SIAM}}=&-t_b\sum_{j=1}^{N_b-1}\sum_\sigma
 \left(c_{j\sigma}^{\dagger}c_{j+1,\sigma}+\mathrm{H.c.}\right)
 \nonumber\\
 &+U\left(n_{d\uparrow}-\frac12\right)
       \left(n_{d\downarrow}-\frac12\right)
 \nonumber\\
 &-V\sum_\sigma
 \left(d_\sigma^{\dagger}c_{1\sigma}+\mathrm{H.c.}\right).
 \label{eq:siam_hamiltonian}
\end{align}
We use $N_b=4$, $t_b=1$, $U=1.5$, and $V=1/\sqrt2$.  Spin-blocked
Jordan--Wigner order gives ten qubits.  The sector
$(N_\uparrow,N_\downarrow)=(3,2)$ has dimension
\begin{equation}
 \dim\mathcal H_{(3,2)}=\binom{5}{3}\binom{5}{2}=100.
 \label{eq:siam_sector_dimension}
\end{equation}
This is one finite-box transfer example,
not a scaling study and not a thermodynamic Kondo-cloud claim.

\begin{widetext}
\inlinetablecaption{tab:benchmark_inventory}{Benchmark inventory.  Every
cell targets the six lowest excited roots in the exact-seed sector.  Pool
sizes refer to the production ordered pools, and exact excited states enter
only post-run audits.}
\begin{center}
\small
\begin{tabular*}{\textwidth}{@{\extracolsep{\fill}}p{0.18\textwidth}p{0.28\textwidth}p{0.15\textwidth}p{0.13\textwidth}p{0.18\textwidth}@{}}
\toprule
Family & cases & encoded register & pool size & seed used here \\
\midrule
Hubbard--Holstein dimer & six $(U,g)$ regimes at $\phmax=1,3$ & 6 or 8 qubits & 169 or 1403 & exact same-sector ground state \\
Peierls--Hubbard dimer & weak--weak, weak--strong, and $U=8$--weak & 6 qubits; sector dimension 16 & 53 & exact same-sector ground state \\
finite-box SIAM & $N_b=4$, $U=1.5$, $V=1/\sqrt2$ & 10 qubits; sector dimension 100 & 204 & exact same-sector ground state \\
\bottomrule
\end{tabular*}
\end{center}
\end{widetext}

\section{Static spectral-window construction}
\label{sec:static_method}

\subsection{Projected record images and the QSE pencil}
\label{sec:qse_pencil}

Let $\Pi_{\mathfrak s}$ project onto the declared conserved sector and define
\begin{equation}
 Q_{\mathfrak s,0}=\Pi_{\mathfrak s}-\ket{\psi_0}\!\bra{\psi_0}.
 \label{eq:q0_projector}
\end{equation}
The implemented pools preserve the sector, so explicit application reduces
to the equivalent reference projection $Q_0=I-\ket{\psi_0}\bra{\psi_0}$.
For an exact seed, a bare power of the Hamiltonian contributes nothing:
\begin{equation}
 Q_{\mathfrak s,0}H_0^k\ket{\psi_0}=0,
 \qquad k\geq0.
 \label{eq:ground_power_vanishes}
\end{equation}
Hamiltonian-power directions are therefore not advertised as a populated
production family.  A nontrivial Krylov route requires a source outside the
ground-state ray, as specified in Sec.~\ref{sec:comparators}.

Let
\begin{equation}
 \mathcal P=\{\Omega_A:A=1,\ldots,N_{\mathcal P}\}
 \label{eq:declared_pool}
\end{equation}
be an ordered deterministic pool, and let $\mathcal B\subseteq\mathcal P$ be
an admitted support.  Its raw projected record images are
\begin{equation}
 \ket{\phi_A}=Q_{\mathfrak s,0}\Omega_A\ket{\psi_0},
 \qquad A\in\mathcal B.
 \label{eq:record_image}
\end{equation}
Define the column map $\Phi_{\mathcal B}e_A=\ket{\phi_A}$ and matrices
\begin{equation}
 S_{\mathcal B}=\Phi_{\mathcal B}^{\dagger}\Phi_{\mathcal B},
 \qquad
 M_{\mathcal B}=\Phi_{\mathcal B}^{\dagger}H_0\Phi_{\mathcal B}.
 \label{eq:qse_matrices}
\end{equation}
The QSE roots and states satisfy
\begin{align}
 M_{\mathcal B}c_r
 &=\varepsilon_r(\mathcal B)S_{\mathcal B}c_r,
 &c_r^{\dagger}S_{\mathcal B}c_s&=\delta_{rs},
 \label{eq:qse_gep}\\
 \ket{\Psi_r(\mathcal B)}
 &=\Phi_{\mathcal B}c_r
 =\sum_{A\in\mathcal B}(c_r)_A\ket{\phi_A}.
 \label{eq:qse_ritz_state}
\end{align}
This is part of the method, not a result.  The main text retains the physical
pencil because it defines both acquisition and all comparators; numerical
whitening details are expanded below and in Appendix~\ref{app:supported_solve}.
Roman $S_{\mathcal B}$ denotes the overlap matrix, whereas calligraphic
$\score_\ell(A)$ will denote the acquisition score.

\subsection{Supported-overlap solution}
\label{sec:supported_overlap}

The solved subspace is the numerically supported image of
$\Phi_{\mathcal B}$ rather than the formal span of every supplied record.  Let
\begin{equation}
 S_{\mathcal B}=U\,\operatorname{diag}(s_1,\ldots,s_{|\mathcal B|})U^\dagger,
 \qquad s_{\max}=\max_i s_i,
 \label{eq:overlap_eigendecomposition}
\end{equation}
and retain
\begin{equation}
 \mathcal I_{\mathcal B,\tau}
 =\left\{i:s_i>\max(\tau_{\rm abs},\tau_{\rm rel}s_{\max})\right\}.
 \label{eq:overlap_support}
\end{equation}
For OS-QSE, direct Davidson, and fixed QSE,
$\tau_{\rm abs}=10^{-12}$ and $\tau_{\rm rel}=10^{-10}$.  The whitening map
and ordinary Hermitian problem are
\begin{align}
 X_{\mathcal B,\tau}
 &=U_{\mathcal I}\operatorname{diag}
   \bigl(s_i^{-1/2}:i\in\mathcal I\bigr),
 \label{eq:supported_whitening}\\
 \widetilde M_{\mathcal B}
 &=X_{\mathcal B,\tau}^{\dagger}
   M_{\mathcal B}X_{\mathcal B,\tau},
 \nonumber\\
 \widetilde M_{\mathcal B}y_r&=\varepsilon_r y_r,
 &c_r&=X_{\mathcal B,\tau}y_r.
 \label{eq:supported_eigenproblem}
\end{align}
The retained rank is
\begin{equation}
 q(\mathcal B)=|\mathcal I_{\mathcal B,\tau}|.
 \label{eq:retained_rank}
\end{equation}
No ridge-shifted pencil is used.  Small overlap directions are removed, and
all claims attach to this supported subspace.  This avoids adding a numerical
parameter that would change the variational problem while preserving the
stability protection that the nonorthogonal solve actually requires
\cite{Epperly2022QSD,Lee2024SamplingQKSD,Getelina2024HardwareNoiseQSE}.

\subsection{Explicit operator pools}
\label{sec:operator_pools}

A ``completion pool'' is an explicit construction, not a proof of
full-sector completeness.  The relevant object is the image set
\begin{equation}
 \mathcal W(\mathcal P,\psi_0)
 =\{Q_{\mathfrak s,0}\Omega_A\ket{\psi_0}:A\in\mathcal P\}.
 \label{eq:pool_image_set}
\end{equation}
Distinct records can become identical or vanish after acting on a particular
seed.

For Hubbard--Holstein, the production pool concatenates the original
physical macro-record pool with an encoded-boson Pauli completion.  If $n_b$
is the number of boson qubits, define
\begin{align}
 \mathcal P_b^{\rm enc}
 &=\{I,X,Y,Z\}^{\otimes n_b}
   \setminus\{I^{\otimes n_b}\},
 \nonumber\\
 \mathcal Z_f&=\{I,Z_1,Z_2,Z_3,Z_4\},
 \label{eq:hh_completion_factors}\\
 \mathcal P_{\rm comp}^{\rm HH}
 &=\{P_b\otimes Z_f:\
     P_b\in\mathcal P_b^{\rm enc},\
     Z_f\in\mathcal Z_f\},
 \label{eq:hh_completion_pool}\\
 \mathcal P_{\rm HH}
 &=\mathcal P_{\rm macro}^{\rm HH}\mathbin{\Vert}
   \mathcal P_{\rm comp}^{\rm HH}.
 \label{eq:hh_production_pool}
\end{align}
The symbol $\Vert$ denotes ordered concatenation, so the macro pool remains
an unchanged prefix.  The completion is arbitrary on the encoded boson
register but uses only identity or one-body $Z$ factors on the four fermion
qubits and therefore does not flip their occupation sector.  It contributes
75 records at $\phmax=1$ and 1275 at $\phmax=3$, giving total pool sizes
169 and 1403.  It is a boson-complete Pauli factor under a deliberately
restricted fermion factor set, not a complete operator basis on the sector.

The Peierls pool is a fixed 53-record macro collection containing
number-conserving fermion hopping, current, and double-excitation records;
bond displacement, momentum, and number records; and their declared mixed
products.  The SIAM calculation uses a 204-record pool formed from physical
impurity/bath blocks together with generalized number, hopping, current, and
paired spin-up--spin-down products.  Appendix~\ref{app:pools} records the
families and their role more fully.  Each route within a comparison receives
the same Hamiltonian and an immutable ordered pool digest.

\subsection{Ideal minimum-cost support}
\label{sec:ideal_construction}

Let the additive structural cost of a support be
\begin{equation}
 K(\mathcal B)=\sum_{A\in\mathcal B}K(A),
 \label{eq:support_cost}
\end{equation}
where $K(A)$ is the deterministic Marrakesh graph-span two-qubit coordinate
of record $A$.  The ideal task is
\begin{equation}
 \mathcal B^\star(\epsilon_E)\in
 \argmin_{\mathcal B\subseteq\mathcal P}K(\mathcal B)
 \label{eq:ideal_support_argmin}
\end{equation}
subject to
\begin{equation}
 q(\mathcal B)\geq R,
 \qquad
 \max_{r\in\mathcalR}
 |\varepsilon_r(\mathcal B)-\varepsilon_r^\star|\leq\epsilon_E.
 \label{eq:ideal_support_constraints}
\end{equation}
This combinatorial statement is useful even though the implementation is
greedy: it specifies the intended method identity.  It is not available
online because the exact energies $\varepsilon_r^\star$ are withheld.

For a supported space $V$ with dimension at least $R$, define the Ky--Fan
functional
\begin{equation}
 \mathcal F_{\mathcalR}(V)
 =\sum_{r\in\mathcalR}\varepsilon_r(V)
 =\min_{\substack{Y^\dagger Y=I_R\\\operatorname{Ran}Y\subseteq V}}
 \Tr(Y^\dagger H_0Y).
 \label{eq:ky_fan_variational}
\end{equation}
The selector evaluates the decrease of this variational functional under a
one-record expansion.  Each proposed-pencil gain is exact for that proposed
supported subspace, but greedy acquisition is not claimed to solve
Eq.~\eqref{eq:ideal_support_argmin} globally.

\subsection{Retained-frame novelty and admissibility}
\label{sec:retained_novelty}

At acquisition round $\ell$, let $\mathcal B_\ell$ be the admitted support
and let
\begin{equation}
 V_{\ell,\tau}
 =\spanop\{\ket{v_1^{(\ell)}},\ldots,\ket{v_{q_\ell}^{(\ell)}}\},
 \qquad
 P_{\ell,\tau}=\sum_{j=1}^{q_\ell}
 \ket{v_j^{(\ell)}}\!\bra{v_j^{(\ell)}}
 \label{eq:retained_frame}
\end{equation}
be the orthonormal frame reconstructed from the retained overlap directions.
Novelty is measured against this frame, not the exact algebraic span of all
admitted images.  A vector discarded by the QSE cutoff must not count as
represented during later acquisition.

For an unused record $A$, write
\begin{equation}
 \ket{w_A}=Q_{\mathfrak s,0}\Omega_A\ket{\psi_0},
 \qquad
 \ket{\widetilde u_{\ell A}}=(I-P_{\ell,\tau})\ket{w_A}.
 \label{eq:candidate_residualization}
\end{equation}
The second operation is evaluated by two complete Gram--Schmidt projection
passes against the retained frame.  The second pass is not cosmetic: a
single pass allowed near-dependent numerical residue to be mistaken for zero
novelty in an earlier implementation.  Define
\begin{equation}
 N_\ell(A)=
 \frac{\|\widetilde u_{\ell A}\|^2}{\|w_A\|^2}.
 \label{eq:novelty_fraction}
\end{equation}
A record is admissible only if
\begin{equation}
 \|w_A\|^2>\tau_{\rm abs},
 \qquad N_\ell(A)\geq\tau_N,
 \qquad \tau_N=10^{-12}.
 \label{eq:candidate_admissibility}
\end{equation}
For an admissible record, the genuinely new direction is
\begin{equation}
 \ket{u_{\ell A}}=
 \frac{\ket{\widetilde u_{\ell A}}}{\|\widetilde u_{\ell A}\|}.
 \label{eq:normalized_novel_direction}
\end{equation}
Novelty is an admissibility floor only.  The production score contains no
additive novelty reward.

\subsection{Full proposed-pencil Ky--Fan gain}
\label{sec:full_pencil_gain}

In the retained orthonormal frame define
\begin{equation}
 \widehat M_\ell=
 \left[\bra{v_i^{(\ell)}}H_0\ket{v_j^{(\ell)}}\right]_{i,j=1}^{q_\ell}.
 \label{eq:retained_hamiltonian}
\end{equation}
For every admissible candidate, the implementation computes the complete new
Hamiltonian row and diagonal,
\begin{equation}
 g_{\ell A,j}=\bra{v_j^{(\ell)}}H_0\ket{u_{\ell A}},
 \qquad
 h_{11,\ell}(A)=\bra{u_{\ell A}}H_0\ket{u_{\ell A}},
 \label{eq:candidate_hamiltonian_row}
\end{equation}
and solves
\begin{equation}
 \widehat M_{\ell,A}^{+}=
 \begin{pmatrix}
   \widehat M_\ell & g_{\ell A}\\
   g_{\ell A}^{\dagger} & h_{11,\ell}(A)
 \end{pmatrix}.
 \label{eq:full_proposed_pencil}
\end{equation}
Ranking therefore uses the candidate's complete interaction with the
retained frame.  It is neither the cosine of a residual angle nor an isolated
$2\times2$ lowest-root surrogate.

Before $R$ supported roots exist, the ordinary $R$-root sum has missing
entries.  The startup rule first computes a common round-local ceiling
\begin{equation}
 \Lambda_\ell=\max_{A\in\mathcal A_\ell}h_{11,\ell}(A),
 \label{eq:rank_filling_ceiling}
\end{equation}
where $\mathcal A_\ell$ is the admissible unused set.  For a supported space
$V$ of dimension $q$, define
\begin{equation}
 \widetilde{\mathcal F}_{\mathcalR}^{(\ell)}(V)
 =\sum_{r=1}^{\min(R,q)}\varepsilon_r(V)
 +(R-q)_+\Lambda_\ell.
 \label{eq:padded_ky_fan}
\end{equation}
The same $\Lambda_\ell$ is used for the current support and every proposal in
the round.  It encourages rank creation without exact spectral information.
Once $q\geq R$, Eq.~\eqref{eq:padded_ky_fan} becomes the ordinary Ky--Fan
sum.  The nonnegative proposal gain is
\begin{equation}
 \Delta_\ell(A)=\left[
 \widetilde{\mathcal F}_{\mathcalR}^{(\ell)}(V_{\ell,\tau})
 -\widetilde{\mathcal F}_{\mathcalR}^{(\ell)}
 \bigl(V_{\ell,\tau}+\spanop\{\ket{u_{\ell A}}\}\bigr)
 \right]_+.
 \label{eq:proposed_pencil_gain}
\end{equation}
The candidate diagonal $h_{11}$ also enters the stopping rule below.  Removing
it after eliminating a separate Ritz-gain expression would silently change
the online window test.

\subsection{Cost-weighted acquisition score}
\label{sec:acquisition_score}

Let $\mathcal A_0$ be the initially eligible pool after structural filters.
The cost scale is fixed once,
\begin{align}
 K_0&=\max_{A\in\mathcal A_0}K(A),
 \nonumber\\
 K_{\max}&=\begin{cases}
 K_0,&K_0>0,\\
 1,&K_0=0,
 \end{cases}
 &\overline K(A)&=\frac{K(A)}{K_{\max}}.
 \label{eq:fixed_cost_normalization}
\end{align}
It is not recomputed as the unused pool shrinks.  By contrast, the gain scale
is round-local,
\begin{align}
 \Delta_\ell^{\max}
 &=\max_{A\in\mathcal A_\ell}\Delta_\ell(A),
 \nonumber\\
 \overline\Delta_\ell(A)=
 \begin{cases}
 \Delta_\ell(A)/\Delta_\ell^{\max},&\Delta_\ell^{\max}>0,\\
 0,&\Delta_\ell^{\max}=0.
 \end{cases}
 \label{eq:round_gain_normalization}
\end{align}
The production score is
\begin{equation}
 \score_\ell(A)=
 \frac{\overline\Delta_\ell(A)}
 {\max\{\kappa_K,\overline K(A)\}^{\alpha_K}},
 \qquad \alpha_K=1,\quad\kappa_K=0.05.
 \label{eq:production_score}
\end{equation}
Away from the cost floor this has the ordering of $\Delta_\ell(A)/K(A)$.
The floor bounds the discount assigned to zero-cost or exceptionally cheap
records and prevents division by zero; it is not a ridge in the QSE pencil.
The largest-score record is admitted, with the original pool index breaking
exact ties.  The retained frame is rebuilt from all admitted raw images
before the next round.  Residual angle, transition visibility, and additive
novelty utility do not enter Eq.~\eqref{eq:production_score}.

It is useful to distinguish score from utility.  In earlier generic selector
code, a spectral utility could be combined additively with novelty and then
discounted by cost.  Under the locked production policy, the additive
novelty term is zero and the spectral utility is simply the normalized
full-pencil gain.  The final ranking quantity is the cost-weighted score
$\score_\ell$.  Thus there is no additional unpublished utility controlling
the reported route.

\subsection{Pool-relative online stopping and statuses}
\label{sec:online_stopping}

For the lowest $\min(R,q_\ell)$ retained Ritz pairs, define
\begin{align}
 \ket{R_r^{(\ell)}}=
 Q_{\mathfrak s,0}(H_0-\varepsilon_r^{(\ell)})
 \ket{\Psi_r^{(\ell)}},
 \nonumber\\
 \rho_\ell=\max_{r\leq\min(R,q_\ell)}\|R_r^{(\ell)}\|.
 \label{eq:target_residual}
\end{align}
Residual convergence alone is insufficient.  A subspace can contain $R$
exact eigenvectors and have zero residual while remaining blind to a lower
state outside it.  The selector therefore asks whether a remaining novel
pool direction can threaten the top of the requested window.  For
$q_\ell\geq R$, define
\begin{equation}
 \mathcal T_\ell=\left\{
 A\in\mathcal A_\ell:
 h_{11,\ell}(A)<\varepsilon_R^{(\ell)}+\epsilon_{\rm res}
 \right\},
 \qquad \epsilon_{\rm res}=10^{-3}.
 \label{eq:window_pressure_set}
\end{equation}
Online acquisition stops successfully only when
\begin{equation}
 q_\ell\geq R,
 \qquad \rho_\ell<\epsilon_{\rm res},
 \qquad \mathcal T_\ell=\varnothing.
 \label{eq:online_stop}
\end{equation}
This is a pool-relative certificate: no remaining admissible image in the
declared pool threatens the retained window under the diagonal test.  It is
not a proof of completeness in the full sector or in an unenumerated family.
After maintenance, the terminal support is solved again and the same
conditions are reevaluated.  The route returns a structured status,
\begin{equation}
 \begin{aligned}
 &\texttt{TARGET\_WINDOW\_CERTIFIED},\qquad
  \texttt{POOL\_EXHAUSTED},\\
 &\texttt{TARGET\_WINDOW\_NOT\_CERTIFIED}.
 \end{aligned}
 \label{eq:terminal_statuses}
\end{equation}
No status silently invokes Davidson or another eigensolver.

\subsection{Atomic singleton exchange}
\label{sec:atomic_exchange}

The maintenance stage adopts the same joint patch object and notation used in
the companion Paper-II exchange construction \cite{PaperII}.  A transaction is
\begin{equation}
 \mathcal X=(\mathcal B^-,\mathcal B^+),
 \qquad
 \mathcal B'=(\mathcal B\setminus\mathcal B^-)\cup\mathcal B^+.
 \label{eq:exchange_patch}
\end{equation}
The production scope is an atomic singleton swap
$|\mathcal B^-|=|\mathcal B^+|=1$, together with the degenerate pure-deletion
case $|\mathcal B^-|=1$, $|\mathcal B^+|=0$.  Pure insertion is disabled.
This is genuine joint $1\leftrightarrow1$ exchange: deletion and insertion are
evaluated as one support.  It does not claim higher-cardinality deletion
ladders.

Shortlists limit how many patches receive a full solve but never certify a
patch.  Deletion candidates combine small target-root coefficient magnitude
and high record cost.  Insertion candidates use a novelty/residual-capture
pre-screen.  That pre-screen belongs only to maintenance; it is not the
acquisition score in Eq.~\eqref{eq:production_score}.  For every proposed
patch, the complete supported QSE pencil on $\mathcal B'$ is rebuilt.

Let
\begin{align}
 m&=\min\{R,q(\mathcal B),q(\mathcal B')\},\\
 \delta\mathcal F&=\sum_{r=1}^{m}
 \bigl[\varepsilon_r(\mathcal B')-\varepsilon_r(\mathcal B)\bigr],
 &\delta K&=K(\mathcal B')-K(\mathcal B).
 \label{eq:exchange_deltas}
\end{align}
With $\tau_{\mathcal F}=10^{-9}$ and
$\eta_{\mathcal F}=10^{-10}$, the spectral--cost gate is
\begin{equation}
 \bigl(\delta\mathcal F\leq\tau_{\mathcal F}\ \wedge\ \delta K<0\bigr)
 \ \vee\
 \bigl(\delta\mathcal F<-\eta_{\mathcal F}\ \wedge\ \delta K\leq0\bigr).
 \label{eq:exchange_dominance}
\end{equation}
The candidate must also preserve retained rank and may not increase the
supported overlap condition estimate by more than a factor of ten.  A failed
solve or guard rejects the entire patch.  Among admissible patches, the route
orders lexicographically by $(\delta\mathcal F,\delta K,$ enumeration index$)$
and commits exactly one.  Rejection leaves the support bit-identical; no
delete occurs before its paired insertion has passed.

The word ``certified'' is attached to this transaction only.  It means that
the recomputed supported pencil satisfies the declared objective, cost, rank,
and conditioning inequalities.  It does not certify every exact root error,
transition amplitude, global minimum-cost support, or completeness outside
the pool.  Current evidence records evaluated exchange locations but not a
current-policy accepted-patch ledger, so exchange is treated below as a
maintenance capability rather than an established cause of the aggregate
savings.

\section{Comparator constructions}
\label{sec:comparators}

\subsection{Independently initialized direct quantum Davidson}
\label{sec:direct_davidson}

Direct quantum Davidson (DQ) is a separate route, not a completion stage of
OS-QSE.  Its starting frame attempts a route-specific linear-response-family
warm start: 10 records for Hubbard--Holstein at $\phmax=1$, 18 at
$\phmax=3$, one for Peierls--Hubbard, and none for SIAM.  It additionally
attempts $R+1=7$ computational-basis probes from the target sector.  The
probes are ordered deterministically by $(H_{ii},i)$, using the Hamiltonian
diagonal and basis index but no exact excited-state information.  Each seed
image is reference-projected, residualized by two Gram--Schmidt passes, and
admitted only when independent; an attempted warm-start record is therefore
not necessarily an admitted direction.

At each iteration DQ solves the same supported Ritz problem.  For the root
with largest unresolved residual, it forms the diagonal-preconditioned
correction
\begin{align}
 \ket{t_r}=D_\delta(\varepsilon_r)^{-1}\ket{R_r},
 \nonumber\\
 [D_\delta(\varepsilon_r)]_{ii}
 &=\operatorname{sgn}_{+}(H_{ii}-\varepsilon_r)
 \nonumber\\[-0.5ex]
 &\quad\times\max\{|H_{ii}-\varepsilon_r|,10^{-8}\}.
 \label{eq:davidson_correction}
\end{align}
Here $\operatorname{sgn}_{+}(x)=-1$ for $x<0$ and $+1$ for $x\geq0$,
so an exactly zero denominator receives the positive $10^{-8}$ floor.
Roots are attempted in descending residual norm until one correction survives
double reorthogonalization and the $10^{-12}$ independence floor.  The route
requests the ground root and six excited roots, uses a native residual
threshold $10^{-14}$, and otherwise ends at its declared cap or when no novel
correction exists.  The paper-facing crossing can occur before native
termination.  Initial record directions retain their record costs; each
synthesized residual direction is charged as one declared first-order
Hamiltonian step.

\subsection{Single-seed real-time Krylov}
\label{sec:single_rtk}

The single-seed RTK comparator constructs a deterministic complex source in
the target sector.  With pseudorandom coefficients fixed before the matrix,
\begin{align}
 \ket{\widetilde\chi}
 &=Q_{\mathfrak s,0}\sum_z(\xi_z+\ii\eta_z)\ket{z},
 \nonumber\\
 \ket{\chi}&=\frac{\ket{\widetilde\chi}}
 {\|\widetilde\chi\|},
 \label{eq:single_rtk_source}
\end{align}
and
\begin{align}
 \ket{\chi_k}&=e^{-\ii k\Delta tH_0}\ket{\chi},
 \qquad k=0,\ldots,d-1,
 \nonumber\\
 V_d^{\RTK}&=\spanop\{\ket{\chi_0},\ldots,\ket{\chi_{d-1}}\},
 \qquad \Delta t=0.5.
 \label{eq:single_rtk_space}
\end{align}
The time-state overlap is truncated at relative cutoff $10^{-12}$.  The
dimension cap is 100 for Hubbard--Holstein and SIAM and 40 for Peierls.
Every Krylov-state preparation is charged; the common source-state
preparation is excluded under the declared boundary.

\subsection{Six-probe block real-time Krylov}
\label{sec:block_rtk}

The block route takes the first six independent directions from the
lowest-diagonal target-sector computational probes.  Let them be
$\{\ket{\chi_j}\}_{j=1}^{6}$.  Define the propagated block states by
\begin{align}
 \ket{\chi_{j,k}}
 &=e^{-\ii k\Delta tH_0}\ket{\chi_j},
 \qquad
 \substack{j=1,\ldots,6,\\ k=0,\ldots,m-1},
 \nonumber\\[-0.5ex]
 V_m^{\rm block}
 &=\spanop\left(
 \{\ket{\psi_0}\}\cup
 \{\ket{\chi_{j,k}}\}_{j,k}
 \right),
 \qquad \Delta t=0.5.
 \label{eq:block_rtk_space}
\end{align}
The block size, time step, and diagonal ordering are fixed before any exact
audit.  Every source preparation per shell and every required Hamiltonian
action is charged.  Single-seed and block RTK are distinct algorithms; their
success sets are never unioned into a case-wise oracle comparator.

\subsection{Complete fixed-class QSE}
\label{sec:fixed_qse}

The fixed baseline solves the supported generalized pencil once on the
complete declared nonadaptive response class.  It performs no selection,
growth, exchange, or residual synthesis.  Hubbard--Holstein uses the
conventional linear-response class; Peierls and SIAM use their corresponding
fixed classes.  A retained rank below six or an error above the target is a
failed outcome, even if the manifold is inexpensive or reproduces one bright
response channel.

\subsection{Comparator scope beyond the numerical matrix}

Variational quantum deflation, subspace-search VQE, multistate contracted VQE,
and qEOM/q-sc-EOM offer important alternative excited-state architectures
\cite{Higgott2019VQD,Nakanishi2019SSVQE,Parrish2019MCVQE,
Ollitrault2020qEOM,AsthanaKumar2023QScEOM}.  They remain literature
comparators here because the repository does not contain matched,
resource-audited implementations under the present task.  Hamiltonian-power
QKSD is likewise implementation-pending; the RTK routes above must not be
renamed as QKSD.  No numerical superiority claim is made against an
unimplemented method.

\section{Observables, matching, and resource coordinates}
\label{sec:evaluation}

\subsection{Retrospective root assignment and first crossings}

OS-QSE, DQ, and fixed-class QSE report their ordered lowest supported roots.
The RTK routes receive a one-to-one Hungarian assignment to exact target roots
only after their online traces are complete \cite{Kuhn1955Hungarian}.  The
assignment cannot select a source, timestep, block, online stopping point, or
support.

Let $n$ index the prefixes in one route's declared growth order.  For energy
tolerance $\epsilon_E$, the reported prefix is
\begin{equation}
 \begin{aligned}
 n_x^\star(\epsilon_E)=\min\biggl\{n:\;&q_x(n)\geq R,\\[-0.4ex]
 &\max_{r\in\mathcalR}
 |\varepsilon_{x,r}(n)-\varepsilon_r^\star|\leq\epsilon_E
 \biggr\}.
 \end{aligned}
 \label{eq:first_crossing}
\end{equation}
Every error and resource coordinate is taken at the same prefix,
\begin{equation}
 K_x^\star(\epsilon_E)=K_x\!\left(n_x^\star(\epsilon_E)\right).
 \label{eq:crossing_cost}
\end{equation}
This is a post-run benchmark diagnostic and is distinct from the online
statuses in Eq.~\eqref{eq:terminal_statuses}.  A route that does not cross
retains its terminal error and charged expenditure but receives no successful
crossing cost.

The earliest-prefix rule is intentionally not a minimum-resource rule.  When
exchange changes the support cost, a later successful prefix can be cheaper
than the first successful prefix.  Three Hubbard--Holstein selected entries
change relative to a legacy minimum-$K$ summary, as shown in
Table~\ref{tab:crossing_reconciliation}.

\begin{widetext}
\inlinetablecaption{tab:crossing_reconciliation}{Reconciliation of the only
three entries for which a later minimum-$K$ successful support differs from
the publication first prefix.  Entries are admitted dimension $k$ and
graph-span cost $K$.}
\begin{center}
\small
\begin{tabular*}{\textwidth}{@{\extracolsep{\fill}}lccccc@{}}
\toprule
Case & route & target & later minimum-$K$ support & publication first prefix & first-prefix error \\
\midrule
weak--strong, $\nph=1$ & OS-QSE & $10^{-4}$ & $k=15$, $K=24$ & $k=14$, $K=26$ & $1.33\times10^{-15}$ \\
weak--strong, $\nph=1$ & OS-QSE & $10^{-6}$ & $k=15$, $K=24$ & $k=14$, $K=26$ & $1.33\times10^{-15}$ \\
weak--strong, $\nph=3$ & OS-QSE & $10^{-4}$ & $k=46$, $K=428$ & $k=42$, $K=484$ & $7.55\times10^{-5}$ \\
\bottomrule
\end{tabular*}
\end{center}
\end{widetext}

\subsection{Transition amplitudes and broadened spectra}
\label{sec:transition_spectra}

For a probe observable $O$, the transition amplitude and strength of a QSE
root are
\begin{equation}
 t_r^O=\bra{\Psi_r}O\ket{\psi_0},
 \qquad I_r^O=|t_r^O|^2.
 \label{eq:transition_strength}
\end{equation}
For route $x$, let $m_x=\min(R,q_x)$ be the number of available roots used in
its response audit.  The zero-temperature broadened response is
\begin{align}
 A_{O,x}(\omega)
 &=\sum_{r=1}^{m_x}I_{x,r}^O
 L_\eta\!\left(\omega-(\varepsilon_r-E_0)\right),
 \nonumber\\
 L_\eta(x)&=\frac{1}{\pi}\frac{\eta}{x^2+\eta^2}.
 \label{eq:broadened_spectrum}
\end{align}
The retrospective exact spectrum is truncated to the same $m_x$ assigned
roots.  Every successful OS-QSE, DQ, and RTK audit has $m_x=R=6$; the
rank-five fixed Peierls manifold necessarily gives a five-root response audit
and remains a failed route for the six-root energy task.  All spectra use
$\eta=0.05$, 801 grid points, and trapezoidal integration.  The
Hubbard--Holstein grid is
$[0,\max(1,\omega_6^\star+6\eta)]$ and uses
\begin{equation}
 \delta_{L^1}(A)=
 \frac{\int|A(\omega)-A^\star(\omega)|\,d\omega}
 {\max\{\int|A^\star(\omega)|\,d\omega,10^{-15}\}},
 \label{eq:hh_spectral_error}
\end{equation}
with transition strengths below $10^{-8}$ treated as dark.  Peierls, SIAM,
and RTK postprocessing extends the upper limit to the largest exact or
approximate pole plus $0.5$ and leaves the normalized error undefined when
the exact integrated weight is at most $10^{-12}$.  The two policies are not
pooled into one cross-family statistic.  Every present spectrum is evaluated
on the terminal manifold, not at the first energy crossing.

\subsection{Three distinct resource coordinates}
\label{sec:resource_boundary}

The online coordinate $K$ is a deterministic Marrakesh graph-span
two-qubit count on the surviving or admitted support.  It is additive by
construction and enters Eq.~\eqref{eq:production_score}.  It is not a letter
for a state, overlap matrix, or universal circuit cost; $C$ is deliberately
not used for this quantity so the notation matches the companion papers.

Post-run Qiskit values use FakeMarrakesh, optimization level 1, and transpiler
seed 7.  They are sums of independently transpiled direction or action
annotations.  Nonzero Pauli words are preserved while complex coefficients
are replaced by their magnitudes for structural compilation.  The resulting
$N_{2q}$ and summed $D_{2q}$ are circuit-structure annotations.  In
particular, summed depth is not the depth of one executable
linear-combination, state-preparation, or measurement circuit.

Logical estimator queries form a third coordinate.  OS-QSE and fixed QSE use
a qubit-wise-commuting record-pencil cover, whereas DQ and RTK use their
declared ungrouped synthesized-basis accounting.  These counts describe the
implemented route ledgers.  They are not shots, and no finite-shot allocation
model is applied.

Finally, compact support is not cumulative acquisition work.  Candidate
screening, rejected proposals, rejected exchange patches, measurement reuse,
source-state preparation, and transition-observable measurements are not
converted into a single common work currency.  We report the available
screening ledger explicitly rather than hiding it, but the main claim remains
about surviving structure under declared coordinates.

\section{Static spectral benchmarks}
\label{sec:static_results}

\subsection{Overall coverage}

Table~\ref{tab:coverage_summary} gives the first high-level result.  OS-QSE
and independently initialized DQ reach every six-root $10^{-6}$ target.
Single-seed RTK and block RTK are incomplete under their fixed source
policies, while the complete fixed classes fail the six-root task in all
sixteen cases.  The fixed manifolds remain scientifically useful response
baselines; their failure is task-specific rather than a claim of universal
inferiority.

\begin{table}[t]
\caption{Six-root $10^{-6}$ coverage under fixed route definitions.  RTK
success sets are complementary and are not unioned into a method.}
\label{tab:coverage_summary}
\centering
\small
\begin{tabular}{@{}lccc@{}}
\toprule
Route & HH (12) & transfer (4) & total (16) \\
\midrule
OS-QSE & 12 & 4 & 16 \\
direct quantum Davidson & 12 & 4 & 16 \\
single-seed RTK & 6 & 4 & 10 \\
six-probe block RTK & 8 & 3 & 11 \\
complete fixed class & 0 & 0 & 0 \\
\bottomrule
\end{tabular}
\end{table}


\subsection{Systematic Hubbard--Holstein results}
\label{sec:hh_static_results}

OS-QSE reaches the six-root $10^{-6}$ target in all twelve
Hubbard--Holstein cells.  DQ also reaches all twelve, while the complete LR-QSE
manifold reaches none and single-seed RTK reaches six.  No OS-QSE result uses
an internal Davidson continuation.

At $10^{-6}$, OS-QSE has smaller $K$ than DQ in all twelve cells.  Its
post-run Qiskit two-qubit annotation and declared logical-query ledger are
also smaller in all twelve.  For $\phmax=1$, OS-QSE crosses with
$K=14$--$28$, versus $56$--$76$ for DQ.  For $\phmax=3$, it crosses with
$K=218$--$526$, versus $308$--$1396$.  The advantage persists at the larger
cutoff but is not uniform in magnitude.

The looser $10^{-4}$ target exposes the dependence on threshold and growth
granularity.  OS-QSE has no larger $K$ in ten of twelve cells, no larger
Qiskit $N_{2q}$ in eleven, and no larger declared query count in all twelve.
The $K$ exceptions are intermediate--weak $\phmax=3$, where the crossings are
$242$ and $240$, and weak--weak $\phmax=3$, where they are $242$ and $172$.
Only the latter remains an exception under the Qiskit annotation.  These
reversals are retained because they prevent a tolerance-independent claim.

Several $\phmax=1$ DQ routes cross only after overshooting to numerical
precision, whereas OS-QSE ranges from numerical precision to
$6.55\times10^{-7}$.  In the $\phmax=3$ matrix, both routes more commonly
cross in the $10^{-7}$--$10^{-6}$ range.  The comparison is threshold matched,
but not every cost difference can be attributed to candidate ordering alone;
the discrete dimension at which a route first records success also matters.

\begin{widetext}
\inlinetablecaption{tab:hh_first_crossings}{Hubbard--Holstein graph-span cost
$K$ at the first $10^{-6}$ six-root crossing.  For single-seed RTK, NR shows
terminal error and terminal charged cost at its declared cap.  Complete
LR-QSE values are terminal and unattainable for the target.}
\begin{center}
\scriptsize
\begin{tabular*}{\textwidth}{@{\extracolsep{\fill}}lcccc@{}}
\toprule
Case & OS-QSE & direct quantum Davidson & single-seed RTK & complete LR-QSE \\
\midrule
intermediate--strong, $\nph=1$ & 28 & 56 & NR: $1.9\times10^{-1}$ @ 99,000 & U: $1.3$ @ 56 \\
intermediate--weak, $\nph=1$ & 28 & 56 & NR: $2.4\times10^{-1}$ @ 99,000 & U: $1.2$ @ 56 \\
strong--strong, $\nph=1$ & 20 & 76 & NR: $5.2\times10^{-1}$ @ 99,000 & U: $7.5$ @ 56 \\
strong--weak, $\nph=1$ & 14 & 76 & NR: $5.3\times10^{-1}$ @ 99,000 & U: $7.8$ @ 56 \\
weak--strong, $\nph=1$ & 26 & 76 & 3,420 & U: $3.1\times10^{-1}$ @ 56 \\
weak--weak, $\nph=1$ & 24 & 76 & 3,420 & U: $6.3\times10^{-1}$ @ 56 \\
intermediate--strong, $\nph=3$ & 440 & 1,192 & 108,528 & U: $6.8\times10^{-1}$ @ 104 \\
intermediate--weak, $\nph=3$ & 260 & 376 & 112,404 & U: $7.3\times10^{-1}$ @ 104 \\
strong--strong, $\nph=3$ & 286 & 308 & NR: $4.8\times10^{-1}$ @ 336,600 & U: $5.5$ @ 104 \\
strong--weak, $\nph=3$ & 218 & 308 & NR: $5.2\times10^{-1}$ @ 336,600 & U: $6.3$ @ 104 \\
weak--strong, $\nph=3$ & 526 & 1,396 & 104,720 & U: $4.6\times10^{-1}$ @ 104 \\
weak--weak, $\nph=3$ & 382 & 512 & 120,360 & U: $2.2\times10^{-1}$ @ 104 \\
\bottomrule
\end{tabular*}
\end{center}
\end{widetext}

The terminal supported ranks equal the admitted sizes in eleven cells.  The
intermediate--weak $\phmax=3$ route retains rank 28 from support 29, showing
that overlap truncation is active without preventing the window.  Terminal
maximum errors remain below $2.5\times10^{-7}$ throughout the $\phmax=3$
matrix and are smaller at $\phmax=1$.  These terminal facts are not substituted
for the first-crossing quantities in Table~\ref{tab:hh_first_crossings}.

\subsection{Structural resource comparison}
\label{sec:resource_results}

At the $10^{-6}$ first crossing, OS-QSE has smaller summed Qiskit $N_{2q}$
than DQ in every Hubbard--Holstein cell and every transfer case.  The SIAM
comparison is especially clear: OS-QSE is slightly more accurate while using
less than one third of the DQ two-qubit annotation.  Appendix
\ref{app:resource_tables} gives the full $N_{2q}/D_{2q}$ and query matrices.

The acquisition ledger makes the accounting boundary concrete.  Across the
twelve Hubbard--Holstein traces, OS-QSE records 325,682 full proposed-pencil
evaluations and 6,752,551 proposed Hamiltonian row elements.  These are real
screening operations but are not charged into the final-support $K$ or Qiskit
columns.  A smaller final manifold can therefore coexist with significant
upstream search.  The present evidence supports compact construction under
the declared support coordinates, not lower cumulative measurement work.

\subsection{Transfer to Peierls--Hubbard and SIAM}
\label{sec:transfer_results}

OS-QSE reaches the $10^{-6}$ window in all three Peierls cases and the
ten-qubit SIAM case; DQ also reaches all four.  In the Peierls cases, OS-QSE
crosses with $K=154$--$202$ and Qiskit annotations $195$--$244$, compared
with DQ $K=300$--$400$ and annotations $417$--$556$.  The OS-QSE logical
query count is 225 in each case, versus 2340 for DQ.  The fixed Peierls class
retains rank five and therefore cannot represent six roots.

The Peierls DQ results overshoot to approximately numerical precision, while
OS-QSE errors range from $1.32\times10^{-9}$ to $5.99\times10^{-8}$.  The
cost result is valid at the matched threshold, but DQ's coarser growth and
overshoot confound a purely selector-based attribution.

The SIAM cell provides the most accuracy-balanced transfer comparison.
OS-QSE reaches error $5.07\times10^{-7}$ with admitted dimension 64 and
retained rank 44.  DQ reaches error $9.60\times10^{-7}$ at dimension and rank
71.  Their Qiskit annotations are 622 and 2142, and their query counts are
4666 and 46,008.  Because OS-QSE is slightly more accurate, this difference
cannot be explained by a looser achieved tolerance.  The $64/44$
support-to-rank reduction also demonstrates active overlap truncation in the
largest encoded example.  One cell remains a transfer demonstration rather
than systematic SIAM generality.

\begin{widetext}
\inlinetablecaption{tab:transfer_results}{Peierls--Hubbard and SIAM results at
the first six-root $10^{-6}$ crossing.  Resource cells are graph-span $K$;
Qiskit $N_{2q}/D_{2q}$ [logical estimator queries].}
\begin{center}
\scriptsize
\begin{tabular*}{\textwidth}{@{\extracolsep{\fill}}lccccc@{}}
\toprule
Case & route & support/rank & achieved error & $K$ & Qiskit $N_{2q}/D_{2q}$ [queries] \\
\midrule
Peierls weak--weak & OS-QSE & $15/15$ & $1.32\times10^{-9}$ & 186 & $195/195$ [225] \\
 & DQ & $12/12$ & $1.55\times10^{-15}$ & 400 & $556/500$ [2340] \\
Peierls weak--strong & OS-QSE & $15/15$ & $5.99\times10^{-8}$ & 202 & $244/228$ [225] \\
 & DQ & $12/12$ & $1.55\times10^{-15}$ & 300 & $417/375$ [2340] \\
Peierls $U=8$--weak & OS-QSE & $12/12$ & $2.78\times10^{-9}$ & 154 & $229/195$ [225] \\
 & DQ & $12/12$ & $3.11\times10^{-15}$ & 300 & $417/375$ [2340] \\
SIAM $N_b=4$ & OS-QSE & $64/44$ & $5.07\times10^{-7}$ & 622 & $622/622$ [4666] \\
 & DQ & $71/71$ & $9.60\times10^{-7}$ & 2142 & $2142/1386$ [46,008] \\
\bottomrule
\end{tabular*}
\end{center}
\end{widetext}

\begin{widetext}
\begin{center}
\includegraphics[width=0.94\textwidth]{../../output/pdf/paper_iii_static_qse_source_anchor_20260829_v2/generality_six_root_bands.png}
\end{center}
\inlinefigurecaption{fig:generality_bands}{Exact and approximate six-root
bands for the Peierls--Hubbard and SIAM transfer cases.  The visibly
overlapping successful bands should be read with the explicit error and cost
values in Table~\ref{tab:transfer_results}; overlap on the plot is not a
resource comparison.}
\end{widetext}

\subsection{Krylov initialization controls target-window coverage}
\label{sec:krylov_results}

The two RTK routes use the same timestep and one-to-one post-run assignment,
but different source subspaces and charging rules.  Single-seed RTK reaches
ten of sixteen windows: six Hubbard--Holstein cells, all three Peierls cases,
and SIAM.  Its Hubbard--Holstein successes split as two of six at $\phmax=1$
and four of six at $\phmax=3$.  Every failure remains visible at its terminal
cap.

Six-probe block RTK reaches eleven of sixteen: all six Hubbard--Holstein
$\phmax=1$ cells, two of six $\phmax=3$ cells, two of three Peierls cases,
and SIAM.  It repairs four single-seed $\phmax=1$ failures and two strong
$\phmax=3$ failures, while losing four windows reached by the single source.
Five cases are reached by both.  Their retrospective union covers all
sixteen, but that union is not a declared algorithm and is never used in a
resource column.

A post-run source-coverage audit explains all five block failures.  In each,
at least one target root has squared projection below $10^{-20}$ onto the
fixed six-probe source block.  The missing target is $r=6$ for Peierls
weak--weak, $r=5$ for intermediate--strong and intermediate--weak
$\phmax=3$, and $r=4$ for weak--strong and weak--weak $\phmax=3$.  A longer
time sequence from the same source block cannot create an eigendirection
absent from that block's spectral support.  Inside an exactly degenerate
eigenspace, individual eigenvector weights are basis dependent; the audit is
therefore interpreted as represented subspace weight, not an invariant label
on an arbitrary eigenvector.

\begin{table*}[t]
\caption{Fixed-policy RTK coverage at the
$10^{-6}$ task.  S and B denote single-seed and block success; a dash denotes
failure within the declared cap.  The final column gives the target index
whose block-source weight is below $10^{-20}$ when applicable.}
\label{tab:krylov_coverage}
\begin{center}
\small
\begin{tabular*}{\textwidth}{@{\extracolsep{\fill}}lcccc@{}}
\toprule
Family/case group & single RTK & block RTK & both & missing block target \\
\midrule
HH $\phmax=1$, weak--weak / weak--strong & S & B & yes & --- \\
HH $\phmax=1$, remaining four regimes & --- & B & no & --- \\
HH $\phmax=3$, strong--weak / strong--strong & --- & B & no & --- \\
HH $\phmax=3$, intermediate--weak / intermediate--strong & S & --- & no & $r=5$ \\
HH $\phmax=3$, weak--weak / weak--strong & S & --- & no & $r=4$ \\
Peierls weak--strong / $U=8$--weak & S & B & yes & --- \\
Peierls weak--weak & S & --- & no & $r=6$ \\
SIAM $N_b=4$ & S & B & yes & --- \\
\bottomrule
\end{tabular*}
\end{center}
\end{table*}

These results make initialization part of the Krylov method definition.  They
do not show that Krylov methods in general fail.  A larger deterministic
block, an adaptive source, or Hamiltonian-power QKSD from a non-eigenstate
source could change both coverage and cost and would be a separately named,
separately charged route.

\subsection{Pool expressivity under the current selector}
\label{sec:pool_ablation_results}

The pool-profile ablation holds the selector, exact seed, overlap cutoff,
dimension cap, stopping rule, and no-fallback boundary fixed.  At the
$10^{-4}$ target, the base physical-macro pool reaches ten of twelve
Hubbard--Holstein cells, whereas encoded-boson completion reaches all twelve.
The repaired cases are weak--strong and intermediate--strong at
$\phmax=3$.  Their completion-pool first crossings occur at $K=484$ and
$428$; the base pool terminates with errors $1.8\times10^{-2}$ and
$1.2\times10^{-2}$ despite much larger terminal expenditure.

Completion-family records appear in every terminal support produced from the
enlarged pool, with counts ranging from three to thirty-six.  The new family
is exercised rather than dormant.  Among the ten cells already solved by the
base pool, the resource effect is not monotone: completion lowers $K$ in some,
leaves it unchanged in others, and raises it in several.  The pool repairs
missing directions and changes the candidate set on which economic selection
acts; it is not a universal cost discount.

\begin{table}[t]
\caption{Controlled Hubbard--Holstein pool-profile result at the six-root
$10^{-4}$ target.  All other route ingredients are fixed.}
\label{tab:pool_ablation_results}
\centering
\small
\begin{tabular}{@{}lcc@{}}
\toprule
Pool & reached & coverage \\
\midrule
physical macro pool & 10/12 & two $\phmax=3$ failures \\
\shortstack[l]{macro + encoded-boson\\completion} & 12/12 & all cells \\
\bottomrule
\end{tabular}
\end{table}

This ablation establishes that the completion repairs two failures under the
current selector.  It does not decompose the total OS-QSE advantage into
independent pool and selection fractions.  A complete same-completion-pool
QSE control, or another matched same-pool selection control, is required for
that stronger causal attribution.

\subsection{Static spectra independently audit the selected states}
\label{sec:static_spectra_results}

Worst-root energy and observable response probe different properties.  For
the twelve Hubbard--Holstein cells, density and displacement channels give a
median normalized spectral $L^1$ error $1.73\times10^{-9}$ for OS-QSE and a
maximum $2.00\times10^{-5}$.  Complete LR-QSE gives median
$3.61\times10^{-3}$ and maximum $1.22\times10^{1}$.  In the two completion
stress cases, the selected broadened peaks and weights track the exact
response precisely where the base physical pool misses the six-root window.

The converse is equally informative.  Complete LR-QSE reproduces the reported
spectra to numerical precision in weak--strong and weak--weak $\phmax=1$
despite failing the six-root energy task.  A root with negligible transition
weight can dominate the worst-root error while barely affecting a chosen
spectrum.  Root-window coverage and response recovery are therefore
complementary diagnostics, not substitutes.

\begin{widetext}
\begin{center}
\includegraphics[width=0.94\textwidth]{../../output/pdf/paper_iii_static_qse_source_anchor_20260829_v2/static_hard_case_spectra.png}
\end{center}
\inlinefigurecaption{fig:stress_spectra}{Exact, OS-QSE, and complete LR-QSE
density and displacement spectra in the two Hubbard--Holstein
completion-pool stress cells.  Every curve is evaluated on the terminal
manifold; first-crossing costs are not assigned to these terminal spectra.}
\end{widetext}

\begin{widetext}
\begin{center}
\includegraphics[width=0.94\textwidth]{../../output/pdf/paper_iii_static_qse_source_anchor_20260829_v2/static_spectral_errors.png}
\end{center}
\inlinefigurecaption{fig:spectral_errors}{Normalized terminal broadened-spectrum
errors for site-0 density and phonon displacement across the twelve
Hubbard--Holstein cells.  Bright-spectrum recovery can occur even when a
six-root energy target is missed, which is why both diagnostics are retained.}
\end{widetext}

For Peierls--Hubbard, OS-QSE normalized errors range from
$4.1\times10^{-14}$ to $4.8\times10^{-12}$ in bond displacement and from
$3.5\times10^{-10}$ to $5.9\times10^{-9}$ in site density where the exact
reference has nonzero weight.  The $U=8$ weak-bond density reference has
numerically zero integrated weight, so its normalized error is undefined.
For SIAM, selected double-occupancy and magnetization errors are
$5.0\times10^{-7}$ and $5.2\times10^{-7}$.  DQ and successful single-seed
RTK also reproduce these transfer spectra accurately; fixed classes retain
visible response errors.

\section{Discussion}
\label{sec:discussion}

\subsection{What the benchmark establishes}

Under the declared six-root task, exact same-cutoff seed convention, pools,
and resource definitions, OS-QSE constructs a successful window in all
sixteen cases.  At the route-local $10^{-6}$ crossing, its graph-span cost,
Qiskit two-qubit structural annotation, and declared query count are each
smaller than DQ in all sixteen.  Hubbard--Holstein supplies the systematic
coupling/cutoff matrix, while Peierls and SIAM test transfer across different
interaction structures and encodings.

The strongest inference is task specific.  Full proposed-pencil Ky--Fan gain
evaluates the quantity optimized by the six-root benchmark, while $K$ biases
growth toward economical directions.  The SIAM result is especially
informative because OS-QSE is slightly more accurate than DQ at the charged
prefix.  The aggregate evidence is therefore consistent with the intended
cost-weighted acquisition mechanism.

``Consistent'' remains important.  DQ overshoots in several cells.  The
completion pool changes both expressivity and the choices presented to the
selector.  Exchange exists as a guarded maintenance operation, but the
current source anchor counts evaluated locations rather than accepted
patches.  The data support the complete declared OS-QSE route; they do not
yet assign a numerical percentage of its savings to cost discounting, pool
completion, or exchange.

\subsection{Why energy windows and response spectra must both remain}

Worst-root error asks whether every requested Ritz value is represented.
A response spectrum weights roots by a particular observable.  OS-QSE is
accurate under both diagnostics for the tested channels, but the LR-QSE
exceptions show they are not equivalent.  A fixed class can be inadequate
for the six-root window yet accurate for a bright response when its missing
roots are dark.  Conversely, six accurate energies do not prove every
transition amplitude, especially near a degeneracy or for an untested probe.
The conclusion is therefore restricted to the declared window and density,
displacement, double-occupancy, and magnetization channels.

\subsection{What the Krylov comparison says}

The single- and block-source matrices identify source coverage as
load-bearing.  Their complementary failures and the five zero-weight audits
show that increasing the number of time samples cannot repair a direction
absent from the fixed source subspace.  A new source block could repair it,
but that changes the algorithm and must carry its own preparation cost.
Accordingly, the evidence establishes complete OS-QSE coverage and lower
declared structural costs relative to the two implemented RTK policies; it
does not establish superiority over Krylov methods as a class.

\subsection{Expressivity and economy are coupled}

The base-pool failures show that economic selection cannot recover a direction
that is absent from the candidate images.  The nonmonotone cost changes after
completion show the converse: a larger pool does not automatically make a
successful cell cheaper.  The observed route relies on a coupled design in
which the pool exposes mixed fermion--boson directions and the selector chooses
within that set.  A same-completion-pool exhaustive control remains the clean
way to separate those two effects.

\subsection{Structural resources are not total work}

The structural claim is intentionally narrower than an end-to-end hardware
claim.  No table includes a complete measurement schedule, shot allocation,
state preparation for an arbitrary linear combination, mitigation, or
wall-clock runtime.  Route-specific query conventions cannot be converted
into a universal shot statement.  The large candidate-screening ledger shows
why final-support compactness and total construction effort must not be
conflated.  A future cumulative ledger should charge screened and rejected
records, accepted and rejected patches, reused matrix elements, and observable
measurements separately.

\subsection{Limitations and remaining mechanism tests}
\label{sec:limitations}

The exact ground-state seed isolates the excited-state problem but does not
demonstrate integration with an approximate variational ground state.  The
bosonic comparisons are same-cutoff and do not establish physical cutoff
convergence.  Finite-shot sampling, mitigation, uncertainty propagation,
runtime, and time dynamics are outside scope.  The single SIAM instance is a
transfer case, not a scaling law.

Three mechanism ledgers remain open.  A current-policy cost-discount-off
ablation is needed to isolate the marginal role of $K$.  An exchange-off
ablation and accepted-patch receipt are needed before exchange can be called a
demonstrated source of savings.  A common cumulative acquisition-work ledger
is needed before structural compactness can imply lower total construction
work.  These tests should hold the pool, seed, supported-overlap cutoff,
stopping rule, and dimension cap fixed.

QKSD remains implementation-pending.  qEOM/q-sc-EOM, VQD, and SSVQE remain
related-work methods rather than matched numerical comparators.  These
absences bound the claims but do not invalidate the sixteen-case comparison
among the implemented routes.

\section{Conclusion}
\label{sec:conclusion}

Cost-weighted full-pencil operator selection constructs accurate six-root QSE
manifolds across the tested Hubbard--Holstein, Peierls--Hubbard, and SIAM
problems.  The selected route reaches all sixteen $10^{-6}$ windows without an
implicit method change and has lower declared structural cost than
independently initialized DQ at the matched crossings.  Transition spectra
show that the resulting manifolds also recover the audited response channels.
The pool and Krylov controls reveal the common mechanism behind the failures:
the target window can only be recovered when the declared candidate or source
subspace exposes its spectral directions.  The remaining work concerns causal
decomposition and cumulative accounting, not repair of the demonstrated
static six-root result.

\appendix

\section{Operator-pool construction and image semantics}
\label{app:pools}

The pool is versioned as an ordered list because ordering resolves exact score
ties and makes deterministic replay meaningful.  Every record also has a
post-seed image.  A family can be algebraically rich yet contribute no novel
direction to a particular seed; conversely, one low-cost Pauli record can
expose a direction absent from a physically named macro family.

\subsection{Hubbard--Holstein macro prefix}

The physical macro prefix contains identity, lifted number-conserving
fermionic single and double directions, site-phonon ladder/displacement and
momentum structures, hopping/current layers, density differences, and mixed
fermion--phonon products.  Records remain at macro granularity: a polynomial
is not silently split into Pauli children during selection.  This is necessary
because $K(A)$ is attached to a record, and mixing macro records with
pre-split children would make ``one admission'' an inconsistent resource
object.

The encoded completion in Eq.~\eqref{eq:hh_completion_pool} appends all
nonidentity boson Pauli words times five sector-preserving fermion factors.
For $n_b=2$ boson qubits, for example, the completion count is
$5(4^2-1)=75$; for $n_b=4$, it is $5(4^4-1)=1275$.  The counts explain the
production pool sizes exactly.  Completion records are not interpreted as
separate physical excitations.  Their purpose is to make the finite encoded
boson algebra visible to the selector while retaining a restricted fermion
factor.

\subsection{Peierls and SIAM pools}

The Peierls pool adapts the physical macro idea to a bond-coupled phonon.  In
addition to fermionic hopping/current and double directions, it contains
$X_b$, its conjugate momentum, phonon number and powers, and products such as
$TX_b$ and current--momentum directions.  Its fixed class contains the
identity, fermionic singles, and bond-phonon ladder directions; its retained
rank five in the reported cases is below the six-root task.

The SIAM pool starts from local impurity/bath physical blocks and adds
generalized number, hopping, and current quadratures within each spin block.
Paired products join one spin-up and one spin-down factor.  The construction
preserves the $(3,2)$ sector.  The selected support of 64 records retaining
rank 44 demonstrates why pool size, support size, and supported dimension
must be reported separately.

\subsection{Seed dependence}

Equation~\eqref{eq:ground_power_vanishes} is a useful warning against generic
pool catalogues.  A bare $H_0^k$ record acting on an exact eigenstate is
annihilated by $Q_{\mathfrak s,0}$.  Likewise, a tangent-derived family is
undefined without a source circuit whose tangent is specified.  We therefore
describe the populated benchmark pools rather than every family that generic
framework code could support.

\section{Supported-pencil geometry and stability}
\label{app:supported_solve}

\subsection{Coordinate invariance}

The matrices in Eq.~\eqref{eq:qse_matrices} depend on the coordinates chosen
inside the same record image.  If $T$ is nonsingular and
$\Phi_{\mathcal B}\mapsto\Phi_{\mathcal B}T$, then
\begin{equation}
 S_{\mathcal B}\mapsto T^\dagger S_{\mathcal B}T,
 \qquad
 M_{\mathcal B}\mapsto T^\dagger M_{\mathcal B}T.
 \label{eq:pencil_congruence}
\end{equation}
The generalized Ritz values are unchanged when the same supported-subspace
convention is applied.  The implementation uses this invariance as a
basis-change sanity check.  A record label is therefore not itself a physical
state; only its image and the supported span matter.

\subsection{Why truncation replaces a ridge}

A ridge would solve a different pencil, for example
$M c=\varepsilon(S+\lambda I)c$, and make the reported roots depend on an
extra deformation.  Supported truncation instead declares that overlap
directions below Eq.~\eqref{eq:overlap_support} are not numerically retained.
In the retained coordinates, $S$ becomes the identity and the Hermitian
Rayleigh--Ritz problem is explicit.

For an isolated retained Ritz pair, the usual residual relation gives the
scale
\begin{equation}
 \sin\angle(\Psi_r,\Psi_r^\star)
 \lesssim
 \frac{\|R_r\|}{\operatorname{gap}_r},
 \label{eq:residual_gap_scale}
\end{equation}
where the denominator is the nearest relevant spectral separation.  Small
overlap eigenvalues amplify perturbations before this residual is formed;
small spectral gaps amplify eigenvector rotation after it is formed.  These
are separate numerical risks.  Near an exact degeneracy, individual
eigenvectors are not invariant, so principal angles and subspace transition
weight are the appropriate diagnostics \cite{DavisKahan1970}.

\subsection{Nested growth and interlacing}

Before exchange and support truncation, append-only growth gives
$V_{\ell}\subseteq V_{\ell+1}$.  Exact Hermitian Rayleigh--Ritz values then
interlace,
\begin{equation}
 \varepsilon_j^{(\ell+1)}
 \leq\varepsilon_j^{(\ell)}
 \leq\varepsilon_{j+d}^{(\ell+1)},
 \qquad
 d=\dim V_{\ell+1}-\dim V_\ell.
 \label{eq:ritz_interlacing}
\end{equation}
An apparent regression can therefore arise from a changed retained support,
root rematching, a changed sector, or finite numerical error rather than from
exact nested Ritz growth.  This is why rank, overlap support, and one-to-one
matching accompany every benchmark trace.

\section{Production selector in operational order}
\label{app:selector_contract}

The equations in the main text can be summarized by the following operational
contract.  It is included to make the manuscript reconstructible without
turning implementation branches into scientific notation.

\begin{enumerate}
\item Construct the ordered pool, per-record $K(A)$, projected raw images,
and immutable pool digest.  Apply the absolute supported-image and structural
filters.  Fix $K_{\max}$ over this initial eligible set.
\item Rebuild the orthonormal retained frame from the current selected raw
images using the same relative overlap cutoff as the final QSE solve.
\item For every unused record, perform two projection passes against the
retained frame, compute $N_\ell(A)$, and remove candidates below the novelty
floor.
\item Apply $H_0$ to every admissible normalized new direction.  Retain the
full Hamiltonian row and $h_{11,\ell}(A)$.
\item If fewer than six supported directions exist, compute the common
rank-filling ceiling $\Lambda_\ell$.  Solve every full proposed matrix in
Eq.~\eqref{eq:full_proposed_pencil} and evaluate $\Delta_\ell(A)$.
\item Normalize gains within the round, divide by the fixed normalized cost
with floor $0.05$, and admit the maximum-$\score_\ell$ record.  Break exact
ties by original pool index.
\item Once six roots exist, evaluate the maximum full-space residual and the
$h_{11}$ window-pressure set.  Stop only when both tests pass.
\item At declared maintenance locations, evaluate atomic singleton patches.
Rejected patches leave the current support unchanged.  Re-solve and repeat
the residual/window certificate on the terminal support.
\item Emit one structured terminal status.  Never change to Davidson inside
the route.
\end{enumerate}

Three implementation details are scientifically load-bearing.  The
Gram--Schmidt pass is repeated.  Novelty is measured against the retained
frame rather than the exact span.  The candidate $h_{11}$ values survive
through the stopping test.  Removing any of these changes the selected trace
or can create false convergence.

\section{Exchange scope and relation to broader add--remove methods}
\label{app:exchange_scope}

Restarted and thick-restart Krylov methods, selected-configuration methods,
and QSE pruning already establish that a useful diagonalization support need
not grow monotonically \cite{WuSimon2000ThickRestart,Sorensen1992IRA}.  The
specific object used here is a measured nonorthogonal operator-pencil
transaction with a cost gate.  Its residual distinction is atomicity: a
delete and insert are scored jointly after re-solving the patched pencil.

The present cardinality is deliberately narrow.  It enumerates singleton
deletions crossed with singleton insertions, plus pure singleton deletions.
A general exchange hierarchy could enumerate complete deletion rungs and
multi-record insertions.  Nothing in the current results demonstrates those
higher-cardinality moves.  Conversely, the singleton restriction does not
reduce the implemented operation to prune-then-append: a joint patch can pass
even when its two halves would not be favorable as sequential independent
moves.

The current numerical authority stores the number of acquisition prefixes at
which exchange maintenance was evaluated.  It does not yet provide a
current-policy aggregate of accepted joint patches.  Older accepted-patch
counts from a different selector policy are not imported.  Accordingly,
exchange is part of the production route and its safety contract, but the
paper does not attribute the observed resource advantage to exchange without
the planned exchange-off ablation.

\section{Detailed structural resource tables}
\label{app:resource_tables}

\begin{table*}[t]
\caption{Post-run structural annotations at the
$10^{-6}$ first crossing.  Numeric cells are Qiskit $N_{2q}/D_{2q}$
[logical queries].  NR cells show terminal-cap resources; fixed LR-QSE is
shown at its complete terminal manifold.}
\label{tab:qiskit_1em6}
\begin{center}
\scriptsize
\begin{tabular*}{\textwidth}{@{\extracolsep{\fill}}lcccc@{}}
\toprule
Case & OS-QSE & direct quantum Davidson & single-seed RTK & complete LR-QSE \\
\midrule
intermediate--strong, $\nph=1$ & $28/28$ [76] & $56/56$ [2448] & NR: $143{,}550/118{,}800$ [90,900] & $56/56$ [121] \\
intermediate--weak, $\nph=1$ & $28/28$ [115] & $56/56$ [2448] & NR: $143{,}550/118{,}800$ [90,900] & $56/56$ [121] \\
strong--strong, $\nph=1$ & $20/20$ [84] & $85/80$ [2448] & NR: $143{,}550/118{,}800$ [90,900] & $56/56$ [121] \\
strong--weak, $\nph=1$ & $14/14$ [28] & $85/80$ [2448] & NR: $143{,}550/118{,}800$ [90,900] & $56/56$ [121] \\
weak--strong, $\nph=1$ & $26/24$ [73] & $85/80$ [2448] & $4959/4104$ [3420] & $56/56$ [121] \\
weak--weak, $\nph=1$ & $24/24$ [126] & $85/80$ [2448] & $4959/4104$ [3420] & $56/56$ [121] \\
intermediate--strong, $\nph=3$ & $461/448$ [1905] & $2248/1192$ [28,896] & $213{,}864/108{,}528$ [52,896] & $104/104$ [845] \\
intermediate--weak, $\nph=3$ & $281/272$ [1398] & $640/376$ [13,920] & $221{,}502/112{,}404$ [54,752] & $104/104$ [845] \\
strong--strong, $\nph=3$ & $307/298$ [1542] & $506/308$ [13,920] & NR: $663{,}300/336{,}600$ [161,600] & $104/104$ [845] \\
strong--weak, $\nph=3$ & $239/236$ [1377] & $506/308$ [13,920] & NR: $663{,}300/336{,}600$ [161,600] & $104/104$ [845] \\
weak--strong, $\nph=3$ & $574/545$ [2001] & $2650/1396$ [33,120] & $206{,}360/104{,}720$ [51,072] & $104/104$ [845] \\
weak--weak, $\nph=3$ & $421/393$ [1767] & $908/512$ [16,896] & $237{,}180/120{,}360$ [58,560] & $104/104$ [845] \\
\bottomrule
\end{tabular*}
\end{center}
\end{table*}

\begin{widetext}
\inlinetablecaption{tab:screening_work}{Recorded Hubbard--Holstein OS-QSE
acquisition activity.  Proposed-pencil evaluations and Hamiltonian row
elements are online screening work; terminal $K$ describes the surviving
support and excludes that screening.  Exchange rungs are evaluated locations,
not accepted-patch counts.}
\begin{center}
\scriptsize
\begin{tabular*}{\textwidth}{@{\extracolsep{\fill}}lrrrrrr@{}}
\toprule
Case & rounds & proposed pencils & row elements & terminal support & terminal $K$ & exchange rungs \\
\midrule
intermediate--strong, $\nph=1$ & 16 & 1,949 & 13,091 & 15 & 28 & 8 \\
intermediate--weak, $\nph=1$ & 16 & 1,948 & 13,087 & 15 & 28 & 8 \\
strong--strong, $\nph=1$ & 15 & 2,115 & 15,364 & 14 & 20 & 7 \\
strong--weak, $\nph=1$ & 15 & 2,115 & 15,364 & 14 & 24 & 7 \\
weak--strong, $\nph=1$ & 16 & 2,142 & 15,553 & 15 & 24 & 8 \\
weak--weak, $\nph=1$ & 16 & 2,092 & 14,889 & 15 & 26 & 8 \\
intermediate--strong, $\nph=3$ & 55 & 74,652 & 2,063,994 & 54 & 452 & 27 \\
intermediate--weak, $\nph=3$ & 30 & 41,423 & 636,892 & 29 & 278 & 15 \\
strong--strong, $\nph=3$ & 32 & 43,970 & 720,432 & 31 & 304 & 16 \\
strong--weak, $\nph=3$ & 24 & 33,166 & 412,504 & 23 & 262 & 12 \\
weak--strong, $\nph=3$ & 55 & 74,621 & 2,062,802 & 54 & 534 & 27 \\
weak--weak, $\nph=3$ & 33 & 45,489 & 768,579 & 32 & 394 & 16 \\
\bottomrule
\end{tabular*}
\end{center}
\end{widetext}

The first table demonstrates that $K$, Qiskit annotation, and query count are
not interchangeable.  The second shows why a final-support resource claim
cannot silently become a total-work claim.  These distinctions remain even
when all coordinates numerically favor the same route at a threshold.

\section{Complete block-Krylov matrix and source audit}
\label{app:block_krylov}

\begin{widetext}
\inlinetablecaption{tab:block_krylov_resources}{Six-probe block-RTK at the
first $10^{-6}$ crossing or terminal cap.  Shape is admitted $k$/retained
rank.  Qiskit is aggregate $N_{2q}/D_{2q}$ over independently prepared block
states.}
\begin{center}
\scriptsize
\begin{tabular*}{\textwidth}{@{\extracolsep{\fill}}lcccccc@{}}
\toprule
Case & status & shape & max error & $K$ & Qiskit $N_{2q}/D_{2q}$ & queries \\
\midrule
intermediate--strong, $\nph=1$ & R & $25/14$ & $1.14\times10^{-13}$ & 720 & $1044/864$ & 5850 \\
intermediate--strong, $\nph=3$ & NR & $97/43$ & $3.71\times10^{-1}$ & 48,960 & $96{,}480/48{,}960$ & 152,096 \\
intermediate--weak, $\nph=1$ & R & $25/14$ & $9.00\times10^{-13}$ & 720 & $1044/864$ & 5850 \\
intermediate--weak, $\nph=3$ & NR & $97/41$ & $4.94\times10^{-1}$ & 48,960 & $96{,}480/48{,}960$ & 152,096 \\
strong--strong, $\nph=1$ & R & $25/14$ & $6.41\times10^{-11}$ & 720 & $1044/864$ & 5850 \\
strong--strong, $\nph=3$ & R & $85/42$ & $7.61\times10^{-7}$ & 37,128 & $73{,}164/37{,}128$ & 116,960 \\
strong--weak, $\nph=1$ & R & $25/14$ & $1.02\times10^{-9}$ & 720 & $1044/864$ & 5850 \\
strong--weak, $\nph=3$ & R & $91/39$ & $4.14\times10^{-7}$ & 42,840 & $84{,}420/42{,}840$ & 133,952 \\
weak--strong, $\nph=1$ & R & $19/13$ & $3.33\times10^{-14}$ & 360 & $522/432$ & 3420 \\
weak--strong, $\nph=3$ & NR & $97/42$ & $5.69\times10^{-1}$ & 48,960 & $96{,}480/48{,}960$ & 152,096 \\
weak--weak, $\nph=1$ & R & $19/13$ & $2.80\times10^{-12}$ & 360 & $522/432$ & 3420 \\
weak--weak, $\nph=3$ & NR & $97/41$ & $2.16\times10^{-1}$ & 48,960 & $96{,}480/48{,}960$ & 152,096 \\
Peierls $U=8$--weak & R & $19/11$ & $1.47\times10^{-14}$ & 1800 & $2502/2250$ & 5700 \\
Peierls weak--strong & R & $19/11$ & $1.39\times10^{-15}$ & 1800 & $2502/2250$ & 5700 \\
Peierls weak--weak & NR & $37/11$ & $3.81\times10^{-1}$ & 9000 & $12{,}510/11{,}250$ & 21,090 \\
SIAM $N_b=4$ & R & $97/88$ & $7.90\times10^{-7}$ & 24,480 & $24{,}480/15{,}840$ & 85,554 \\
\bottomrule
\end{tabular*}
\end{center}
\end{widetext}

\begin{center}
\begin{minipage}{\columnwidth}
\inlinetablecaption{tab:block_source_failures}{Post-run source coverage for
all block-RTK failures.}
\centering
\scriptsize
\setlength{\tabcolsep}{3pt}
\resizebox{\columnwidth}{!}{%
\begin{tabular}{@{}lccc@{}}
\toprule
Case & absent target & source weight & terminal error \\
\midrule
Peierls weak--weak & $r=6$ & $2.7\times10^{-29}$ & $3.81\times10^{-1}$ \\
intermediate--strong, $\nph=3$ & $r=5$ & $1.7\times10^{-26}$ & $3.71\times10^{-1}$ \\
intermediate--weak, $\nph=3$ & $r=5$ & $4.0\times10^{-30}$ & $4.94\times10^{-1}$ \\
weak--strong, $\nph=3$ & $r=4$ & $2.1\times10^{-29}$ & $5.69\times10^{-1}$ \\
weak--weak, $\nph=3$ & $r=4$ & $1.1\times10^{-30}$ & $2.16\times10^{-1}$ \\
\bottomrule
\end{tabular}
}
\end{minipage}
\end{center}

\clearpage
\section{Pool-profile evidence in full}
\label{app:pool_ablation}

\begin{widetext}
\inlinetablecaption{tab:pool_ablation_full}{Current-selector pool comparison at
the first $10^{-4}$ crossing.  R gives crossing dimension $k$, graph-span
$K$, and achieved error; NR gives terminal error and $K$.  The final column
counts completion-family records in the completion-pool terminal support.}
\begin{center}
\scriptsize
\begin{tabular*}{\textwidth}{@{\extracolsep{\fill}}lccc@{}}
\toprule
Case & physical macro pool & completion pool & completion records \\
\midrule
weak--weak, $\nph=1$ & R: $k=12$, $K=18$, $5.3\times10^{-5}$ & R: $k=12$, $K=18$, $5.3\times10^{-5}$ & 4 \\
intermediate--weak, $\nph=1$ & R: $k=12$, $K=52$, $9.7\times10^{-5}$ & R: $k=12$, $K=26$, $6.1\times10^{-5}$ & 5 \\
strong--weak, $\nph=1$ & R: $k=10$, $K=10$, $1.9\times10^{-5}$ & R: $k=10$, $K=14$, $6.5\times10^{-7}$ & 4 \\
weak--strong, $\nph=1$ & R: $k=13$, $K=16$, $3.4\times10^{-14}$ & R: $k=14$, $K=26$, $1.3\times10^{-15}$ & 7 \\
intermediate--strong, $\nph=1$ & R: $k=14$, $K=72$, $2.2\times10^{-5}$ & R: $k=14$, $K=26$, $6.3\times10^{-6}$ & 7 \\
strong--strong, $\nph=1$ & R: $k=12$, $K=18$, $1.7\times10^{-5}$ & R: $k=10$, $K=12$, $1.7\times10^{-5}$ & 3 \\
weak--weak, $\nph=3$ & R: $k=14$, $K=226$, $3.1\times10^{-5}$ & R: $k=14$, $K=242$, $5.3\times10^{-5}$ & 13 \\
intermediate--weak, $\nph=3$ & R: $k=14$, $K=968$, $1.1\times10^{-5}$ & R: $k=18$, $K=242$, $8.4\times10^{-6}$ & 16 \\
strong--weak, $\nph=3$ & R: $k=10$, $K=160$, $4.9\times10^{-6}$ & R: $k=10$, $K=160$, $4.9\times10^{-6}$ & 11 \\
weak--strong, $\nph=3$ & NR: $1.8\times10^{-2}$ at $K=4630$ & R: $k=42$, $K=484$, $7.5\times10^{-5}$ & 35 \\
intermediate--strong, $\nph=3$ & NR: $1.2\times10^{-2}$ at $K=7418$ & R: $k=50$, $K=428$, $6.0\times10^{-6}$ & 36 \\
strong--strong, $\nph=3$ & R: $k=12$, $K=216$, $3.8\times10^{-5}$ & R: $k=12$, $K=216$, $3.8\times10^{-5}$ & 18 \\
\bottomrule
\end{tabular*}
\end{center}
\end{widetext}

The two repaired cells also show why a terminal failure must retain its
expenditure.  Their base-pool costs are not successful alternatives at the
target.  Conversely, the completion pool is not uniformly economical in the
ten jointly successful cells.  These observations are part of the mechanism
claim, not merely supplementary bookkeeping.

\section{Terminal manifolds and channelwise spectral audits}
\label{app:spectral_tables}

\begin{widetext}
\inlinetablecaption{tab:terminal_hh}{Terminal Hubbard--Holstein OS-QSE
manifolds.  The response columns are the mean normalized $L^1$ error over
site-0 density and phonon displacement for OS-QSE and complete LR-QSE.}
\begin{center}
\scriptsize
\begin{tabular*}{\textwidth}{@{\extracolsep{\fill}}lcccccc@{}}
\toprule
Case & support/rank & terminal $K$ & Qiskit $N_{2q}/D$ [queries] & max root error & OS-QSE $L^1$ & LR-QSE $L^1$ \\
\midrule
intermediate--strong, $\nph=1$ & $15/15$ & 28 & $28/28$ [76] & $4.00\times10^{-15}$ & $1.56\times10^{-14}$ & $1.16$ \\
intermediate--weak, $\nph=1$ & $15/15$ & 28 & $28/28$ [115] & $1.17\times10^{-13}$ & $3.15\times10^{-14}$ & $6.10$ \\
strong--strong, $\nph=1$ & $14/14$ & 20 & $20/20$ [84] & $1.66\times10^{-8}$ & $4.01\times10^{-14}$ & $8.40\times10^{-2}$ \\
strong--weak, $\nph=1$ & $14/14$ & 24 & $24/24$ [89] & $8.26\times10^{-10}$ & $3.98\times10^{-13}$ & $4.62\times10^{-1}$ \\
weak--strong, $\nph=1$ & $15/15$ & 24 & $24/24$ [73] & $1.40\times10^{-14}$ & $1.26\times10^{-14}$ & $1.20\times10^{-14}$ \\
weak--weak, $\nph=1$ & $15/15$ & 26 & $26/26$ [128] & $4.12\times10^{-15}$ & $2.08\times10^{-14}$ & $1.57\times10^{-14}$ \\
intermediate--strong, $\nph=3$ & $54/54$ & 452 & $473/460$ [1905] & $1.17\times10^{-7}$ & $3.18\times10^{-7}$ & $2.71\times10^{-3}$ \\
intermediate--weak, $\nph=3$ & $29/28$ & 278 & $299/290$ [1476] & $1.07\times10^{-7}$ & $1.34\times10^{-6}$ & $2.37\times10^{-1}$ \\
strong--strong, $\nph=3$ & $31/31$ & 304 & $325/316$ [1617] & $1.05\times10^{-7}$ & $1.04\times10^{-5}$ & $1.36\times10^{-5}$ \\
strong--weak, $\nph=3$ & $23/23$ & 262 & $283/274$ [1452] & $1.22\times10^{-7}$ & $1.88\times10^{-7}$ & $4.43\times10^{-1}$ \\
weak--strong, $\nph=3$ & $54/54$ & 534 & $582/551$ [2001] & $2.31\times10^{-8}$ & $8.36\times10^{-8}$ & $2.20\times10^{-1}$ \\
weak--weak, $\nph=3$ & $32/32$ & 394 & $436/404$ [1777] & $2.45\times10^{-7}$ & $1.77\times10^{-6}$ & $2.58$ \\
\bottomrule
\end{tabular*}
\end{center}
\end{widetext}

\begin{table*}[t]
\caption{Terminal normalized spectral $L^1$
errors for transfer probes.  The Peierls entries are bond displacement/site-0
density; SIAM entries are double occupancy/magnetization.  ``--'' denotes a
zero-reference-weight channel for which normalized error is undefined.}
\label{tab:transfer_spectra}
\begin{center}
\small
\begin{tabular*}{\textwidth}{@{\extracolsep{\fill}}lcccc@{}}
\toprule
Case & OS-QSE & direct Davidson & single RTK & complete fixed class \\
\midrule
Peierls weak--weak & $1.3\times10^{-13}/5.9\times10^{-9}$ & $3.0\times10^{-15}/1.5\times10^{-14}$ & $6.1\times10^{-15}/7.7\times10^{-15}$ & $2.1\times10^{-2}/2.3$ \\
Peierls weak--strong & $4.8\times10^{-12}/3.5\times10^{-10}$ & $1.2\times10^{-14}/1.2\times10^{-14}$ & $1.8\times10^{-13}/2.7\times10^{-13}$ & $2.9\times10^{-1}/5.2$ \\
Peierls $U=8$--weak & $4.1\times10^{-14}$/-- & $5.9\times10^{-15}$/-- & $2.5\times10^{-14}$/-- & $5.6\times10^{-3}$/-- \\
SIAM $N_b=4$ & $5.0\times10^{-7}/5.2\times10^{-7}$ & $4.4\times10^{-14}/2.7\times10^{-14}$ & $5.6\times10^{-7}/1.1\times10^{-7}$ & $7.5\times10^{-2}/1.0\times10^{-1}$ \\
\bottomrule
\end{tabular*}
\end{center}
\end{table*}

These tables also document that the spectra are terminal audits.  A future
first-crossing spectral study would be a distinct analysis and would need to
charge observable data at that prefix rather than borrowing terminal curves.

\section{Evidence authority, determinism, and data boundary}
\label{app:evidence}

The numerical authority for this draft is the thirteen-page
\sourceanchor, generated from immutable route JSONs and receipts by the
tracked source-anchor builder.  Its
publication-crossing reconstruction, complete resource matrices, failure
statuses, exact-reference identity checks, and spectral conventions supersede
legacy generated manuscript fragments.  The source-anchor PDF SHA-256 at this
revision is
\begin{quote}
\small\ttfamily
4a9d1224785484d34e448617a486bb9\\
12d61da7a93d76272cf77da9853e841a4.
\end{quote}
Two consecutive builds reproduce the PDF byte for byte.  All four
Peierls/SIAM OS-QSE cases also replay with identical selected supports, stop
reasons, and final selector residuals.  The Hubbard--Holstein production
traces retain ordered support indices and online receipt digests.

The source anchor is primarily a numerical contract.  This manuscript is the
scientific framing: it derives the method, explains the comparison boundary,
interprets mechanisms, and states limitations.  Where a numerical value in a
legacy table disagrees with the source anchor, the source anchor controls.
Where prose about the selector disagrees with the locked code, the code-level
policy and its tests control and the anchor must be regenerated.  This rule
was applied to the fixed initial-pool cost normalization in
Eq.~\eqref{eq:fixed_cost_normalization}.

The current repository evidence intentionally leaves several entries open:
cost-discount-off, exchange-off with accepted-patch ledger, common cumulative
acquisition work, Hamiltonian-power QKSD, approximate-seed robustness, and
physical bosonic-cutoff convergence.  An absent entry is not imputed from a
nearby route, and a failed route is not dropped from a denominator.

\begin{thebibliography}{99}

\bibitem{PaperI}
J. Strobel and J. Simoni,
``Geometric and resource-aware adaptive ansatz construction for fermion--boson Hamiltonian systems,''
companion paper (Paper I of this series).

\bibitem{PaperII}
J. Strobel and J. Simoni,
``Append--prune McLachlan dynamics,''
companion paper (Paper II of this series).

\bibitem{Hubbard1963}
J. Hubbard,
``Electron correlations in narrow energy bands,''
Proc. R. Soc. A \textbf{276}, 238--257 (1963).

\bibitem{Holstein1959}
T. Holstein,
``Studies of polaron motion: Parts I and II,''
Ann. Phys. \textbf{8}, 325--342 and 343--389 (1959).

\bibitem{SuSchriefferHeeger1979}
W. P. Su, J. R. Schrieffer, and A. J. Heeger,
``Solitons in polyacetylene,''
Phys. Rev. Lett. \textbf{42}, 1698--1701 (1979).

\bibitem{McClean2017QSE}
J. R. McClean, M. E. Kimchi-Schwartz, J. Carter, and W. A. de Jong,
``Hybrid quantum-classical hierarchy for mitigation of decoherence and determination of excited states,''
Phys. Rev. A \textbf{95}, 042308 (2017).

\bibitem{Colless2018QSE}
J. I. Colless \textit{et al.},
``Computation of molecular spectra on a quantum processor with an error-resilient algorithm,''
Phys. Rev. X \textbf{8}, 011021 (2018).

\bibitem{Urbanek2020VQSE}
M. Urbanek, D. Camps, R. Van Beeumen, and W. A. de Jong,
``Chemistry on quantum computers with virtual quantum subspace expansion,''
J. Chem. Theory Comput. \textbf{16}, 5425--5431 (2020).

\bibitem{Yoshioka2022GQSE}
N. Yoshioka, H. Hakoshima, Y. Matsuzaki, Y. Tokunaga, Y. Suzuki, and S. Endo,
``Generalized quantum subspace expansion,''
Phys. Rev. Lett. \textbf{129}, 020502 (2022).

\bibitem{CortesGray2022QKSD}
C. L. Cortes and S. K. Gray,
``Quantum Krylov subspace algorithms for ground- and excited-state energy estimation,''
Phys. Rev. A \textbf{105}, 022417 (2022).

\bibitem{Motta2024SubspaceReview}
M. Motta \textit{et al.},
``Subspace methods for electronic structure simulations on quantum computers,''
Electron. Struct. \textbf{6}, 013001 (2024).

\bibitem{Jamet2022GreenQSE}
F. Jamet, A. Agarwal, and I. Rungger,
``Quantum subspace expansion algorithm for Green's functions,''
arXiv:2205.00094 (2022).

\bibitem{Ollitrault2020qEOM}
P. J. Ollitrault \textit{et al.},
``Quantum equation of motion for computing molecular excitation energies on a noisy quantum processor,''
Phys. Rev. Research \textbf{2}, 043140 (2020).

\bibitem{Tkachenko2024QDavidson}
N. V. Tkachenko \textit{et al.},
``Quantum Davidson algorithm for excited states,''
Quantum Sci. Technol. \textbf{9}, 035012 (2024).

\bibitem{Epperly2022QSD}
E. N. Epperly, L. Lin, and Y. Nakatsukasa,
``A theory of quantum subspace diagonalization,''
SIAM J. Matrix Anal. Appl. \textbf{43}, 1263--1290 (2022).

\bibitem{Lee2024SamplingQKSD}
G. Lee, D. Lee, and J. Huh,
``Sampling error analysis in quantum Krylov subspace diagonalization,''
Quantum \textbf{8}, 1477 (2024).

\bibitem{Getelina2024HardwareNoiseQSE}
J. C. Getelina, P. Sharma, T. Iadecola, P. P. Orth, and Y.-X. Yao,
``Quantum subspace expansion in the presence of hardware noise,''
APL Quantum \textbf{1}, 036127 (2024).

\bibitem{Kuhn1955Hungarian}
H. W. Kuhn,
``The Hungarian method for the assignment problem,''
Naval Res. Logist. Q. \textbf{2}, 83--97 (1955).

\bibitem{Higgott2019VQD}
O. Higgott, D. Wang, and S. Brierley,
``Variational quantum computation of excited states,''
Quantum \textbf{3}, 156 (2019).

\bibitem{Nakanishi2019SSVQE}
K. M. Nakanishi, K. Mitarai, and K. Fujii,
``Subspace-search variational quantum eigensolver for excited states,''
Phys. Rev. Research \textbf{1}, 033062 (2019).

\bibitem{Parrish2019MCVQE}
R. M. Parrish \textit{et al.},
``Quantum computation of electronic transitions using a variational quantum eigensolver,''
Phys. Rev. Lett. \textbf{122}, 230401 (2019).

\bibitem{AsthanaKumar2023QScEOM}
A. Asthana \textit{et al.},
``Quantum self-consistent equation-of-motion method for computing molecular excitation energies, ionization potentials, and electron affinities on a quantum computer,''
Chem. Sci. \textbf{14}, 2405--2418 (2023).

\bibitem{DavisKahan1970}
C. Davis and W. M. Kahan,
``The rotation of eigenvectors by a perturbation. III,''
SIAM J. Numer. Anal. \textbf{7}, 1--46 (1970).

\bibitem{WuSimon2000ThickRestart}
K. Wu and H. Simon,
``Thick-restart Lanczos method for large symmetric eigenvalue problems,''
SIAM J. Matrix Anal. Appl. \textbf{22}, 602--616 (2000).

\bibitem{Sorensen1992IRA}
D. C. Sorensen,
``Implicit application of polynomial filters in a $k$-step Arnoldi method,''
SIAM J. Matrix Anal. Appl. \textbf{13}, 357--385 (1992).

\end{thebibliography}

\end{document}
